% Generated mechanically from frozen input inputs/p11/ja/examples-defos.tex.
% Do not edit this generated file; edit the bound locale source instead.

% OK, start here.
%

\setcounter{chapter}{92}
\chapter{変形問題}


\phantomsection
\label{examples-defos-section-phantom}


\section{序論}
\label{examples-defos-section-introduction}

\noindent
本章の目的は、『形式的変形理論』、『変形理論』および『余接複体』の
各章で展開した一般理論の例を具体的に検討することである。

\medskip\noindent
Schlessinger による論文 \cite{Sch} の第 3 節でも、
いくつかの例が論じられている。






\section{変形問題の例}
\label{examples-defos-section-examples}

\noindent
ここに入れるべき事項の一覧：
\begin{enumerate}
\item スキームの変形：
\begin{enumerate}
\item Rim--Schlessinger 条件。
\item 接空間の計算。
\item 無限小変形の計算。
\item アフィン超曲面の変形圏。
\end{enumerate}
\item 層の変形（たとえば $X/S$、$S$ の有限型の点 $s$、および
$X_s$ 上の準連接層 $\mathcal{F}_s$ を固定する）。
\item 代数空間の変形（スキームの変形と非常によく似ており、
もしかするとさらに容易かもしれない）。
\item 写像の変形（たとえばスキーム間の射。終域と始域の双方、
またはいずれか一方を固定できる）。
\item ここにさらに追加。
\end{enumerate}





\section{一般的な手順}
\label{examples-defos-section-general}

\noindent
本節では、以下のいくつかの例を論じるための手順を示す。

\medskip\noindent
手順 I. 各節でネーター環 $\Lambda$ と有限環準同型
$\Lambda \to k$ を固定する。ここで $k$ は体である。
通常どおり、基礎圏を
$\mathcal{C}_\Lambda = \mathcal{C}_{\Lambda, k}$
とする。『形式的変形理論』の定義
\ref{formal-defos-definition-CLambda} を参照されたい。

\medskip\noindent
手順 II. 各節で、$\mathcal{C}_\Lambda$ 上のグルーポイドを
ファイバーとする余ファイバー圏 $\mathcal{F}$ を定義する。
場合によっては、その代わりに関手
$F : \mathcal{C}_\Lambda \to \textit{Sets}$ を考える。

\medskip\noindent
手順 III. 『形式的変形理論』の節
\ref{formal-defos-section-RS-condition} で論じた
Rim--Schlessinger 条件 (RS) を $\mathcal{F}$ がどの程度満たすかを説明する。
同様に、$\mathcal{F}$ が (S1) と (S2) をどの程度満たすか、または
$F$ が対応する Schlessinger 条件 (H1) と (H2) をどの程度満たすかを
論じることもある。『形式的変形理論』の節
\ref{formal-defos-section-schlessinger-conditions} を参照されたい。

\medskip\noindent
手順 IV. $x_0$ を $\mathcal{F}(k)$ の対象、言い換えれば $k$ 上の
$\mathcal{F}$ の対象とする。本章では記法
$$
\Deformationcategory_{x_0} = \mathcal{F}_{x_0}
$$
を用いて、『形式的変形理論』の注意
\ref{formal-defos-remark-localize-cofibered-groupoid} で構成した
前変形圏を表す。$\mathcal{F}$ が (RS) を満たすならば、
$\Deformationcategory_{x_0}$ は変形圏であり
（『形式的変形理論』の補題
\ref{formal-defos-lemma-localize-RS}）、
さらに (S1) と (S2) を満たす
（『形式的変形理論』の補題
\ref{formal-defos-lemma-RS-implies-S1-S2}）。
(S1) と (S2) が満たされる場合、重要な問題は接空間
$$
T\Deformationcategory_{x_0} = T_{x_0}\mathcal{F} = T\mathcal{F}_{x_0}
$$
（『形式的変形理論』の注意
\ref{formal-defos-remark-tangent-space-cofibered-groupoid} および定義
\ref{formal-defos-definition-tangent-space} を参照）が有限次元かどうかである。
実際、これにより $\Deformationcategory_{x_0}$ が半普遍的形式対象をもつことが
保証される（『形式的変形理論』の補題
\ref{formal-defos-lemma-versal-object-existence}）。

\medskip\noindent
手順 V. $\mathcal{F}$ が手順 IV を通過したならば、次の問題は、
$x_0$ の無限小自己同型からなる $k$-ベクトル空間
$$
\text{Inf}(\Deformationcategory_{x_0}) = \text{Inf}_{x_0}(\mathcal{F})
$$
が有限次元かどうかである。有限次元であれば、
$\Deformationcategory_{x_0}$ は $\mathcal{C}_\Lambda$ 上の関手からなる
滑らかなプロ表現可能グルーポイドによる表示をもつ。
『形式的変形理論』の定理
\ref{formal-defos-theorem-presentation-deformation-groupoid} を参照されたい。




\section{有限射影加群}
\label{examples-defos-section-finite-projective-modules}

\noindent
本節は準備運動にすぎない。当然ながら、有限射影加群には
「モジュライ」は存在しないはずである。

\begin{example}[有限射影加群]
\label{examples-defos-example-finite-projective-modules}
圏 $\mathcal{F}$ を次のように定義する。
\begin{enumerate}
\item 対象は、$\mathcal{C}_\Lambda$ の対象 $A$ と有限射影
$A$-加群 $M$ からなる対 $(A, M)$ である。
\item 射 $(f, g) : (B, N) \to (A, M)$ は、
$\mathcal{C}_\Lambda$ における射 $f : B \to A$ と、$f$-線形で
同型 $N \otimes_{B, f} A \cong M$ を誘導する写像
$g : N \to M$ からなる。
\end{enumerate}
関手 $p : \mathcal{F} \to \mathcal{C}_\Lambda$ は $(A, M)$ を $A$ に、
$(f, g)$ を $f$ に送る。$p$ がグルーポイドをファイバーとする
余ファイバー圏であることは明らかである。有限次元 $k$-ベクトル空間
$V$ に対し、$x_0 = (k, V)$ を対応する $\mathcal{F}(k)$ の対象とする。
次のようにおく。
$$
\Deformationcategory_V = \mathcal{F}_{x_0}
$$
\end{example}

\noindent
局所環上の有限射影加群はすべて有限自由であるから
（『代数』の補題 \ref{algebra-lemma-finite-projective}）、
次を得る。
$$
\begin{matrix}
\mathcal{F}(A)\text{ の対象の} \\
\text{同型類}
\end{matrix}
= \coprod\nolimits_{n \geq 0} \{*\}
$$
これは $\mathcal{F}$ の変形理論が本質的に自明であることを意味するが、
簡単な例を与えるため、なお節 \ref{examples-defos-section-general} で示した手順を実行する。

\begin{lemma}
\label{examples-defos-lemma-finite-projective-modules-RS}
例 \ref{examples-defos-example-finite-projective-modules} は
Rim--Schlessinger 条件 (RS) を満たす。特に、任意の有限次元
$k$-ベクトル空間 $V$ に対して $\Deformationcategory_V$ は変形圏である。
\end{lemma}

\begin{proof}
$A_1 \to A$ と $A_2 \to A$ を $\mathcal{C}_\Lambda$ の射とし、
$A_2 \to A$ は全射であると仮定する。『形式的変形理論』の補題
\ref{formal-defos-lemma-RS-2-categorical} により、関手
$\mathcal{F}(A_1 \times_A A_2) \to
\mathcal{F}(A_1) \times_{\mathcal{F}(A)} \mathcal{F}(A_2)$
が圏同値であることを示せば十分である。

\medskip\noindent
したがって、$A_1 \times_A A_2$ 上の有限射影加群の圏が、
$A$ 上の有限射影加群の圏を介した、$A_1$ 上および $A_2$ 上の
有限射影加群の圏のファイバー積と同値であることを示す必要がある。
これは『代数詳論』の補題
\ref{more-algebra-lemma-finitely-presented-module-over-fibre-product}
の特別な場合である。逆関手は、$M_1$ が有限射影 $A_1$-加群、
$M_2$ が有限射影 $A_2$-加群であり、
$\varphi : M_1 \otimes_{A_1} A \to M_2 \otimes_{A_2} A$
が $A$-加群の同型である三つ組 $(M_1, M_2, \varphi)$ を、有限射影
$A_1 \times_A A_2$-加群 $M_1 \times_\varphi M_2$ に送ることを思い出そう。
\end{proof}

\begin{lemma}
\label{examples-defos-lemma-finite-projective-modules-TI}
例 \ref{examples-defos-example-finite-projective-modules} において、
$V$ を有限次元 $k$-ベクトル空間とする。このとき
$$
T\Deformationcategory_V = (0)
\;\text{and}\;
\text{Inf}(\Deformationcategory_V) = \text{End}_k(V)
$$
はいずれも有限次元である。
\end{lemma}

\begin{proof}
例 \ref{examples-defos-example-finite-projective-modules} の $\mathcal{F}$ に対して
$x_0 = (k, V) \in \Ob(\mathcal{F}(k))$ とおく。
$T\Deformationcategory_V = T_{x_0}\mathcal{F}$ は、双対数環
$k[\epsilon]$ 上の $\mathcal{F}$ の対象 $x$ と、$k[\epsilon] \to k$
上にある $\mathcal{F}$ の射 $\alpha : x \to x_0$ からなる対
$(x, \alpha)$ の同型類の集合であった。

\medskip\noindent
同型を除けば、$k[\epsilon]$ 上の有限射影加群 $M$ と、同型
$M \otimes_{k[\epsilon]} k \to V$ を誘導する $k[\epsilon]$-線形写像
$\alpha : M \to V$ からなる対 $(M, \alpha)$ は一意に存在する。
たとえば $V = k^{\oplus n}$ ならば、$M = k[\epsilon]^{\oplus n}$ とし、
$\alpha$ には明らかな写像をとる。

\medskip\noindent
同様に、$\text{Inf}(\Deformationcategory_V) = \text{Inf}_{x_0}(\mathcal{F})$
は、$k[\epsilon]$ 上の $x_0$ の自明変形 $x'_0$ の自己同型の集合である。
詳しくは『形式的変形理論』の定義
\ref{formal-defos-definition-infinitesimal-auts} を参照されたい。

\medskip\noindent
第 2 段落の $(M, \alpha)$ をとると、
$\text{Inf}_{x_0}(\mathcal{F})$ の元は
$\gamma \bmod \epsilon = \text{id}$ を満たす自己同型
$\gamma : M \to M$ である。したがって
$\gamma = \text{id}_M + \epsilon \psi$ と書ける。ここで
$\psi : M/\epsilon M \to M/\epsilon M$ は $k$-線形である。
$\alpha$ を用いて $\psi$ を $\text{End}_k(V)$ の元とみなせるので、
証明は完了する。
\end{proof}


\section{群の表現}
\label{examples-defos-section-representations}

\noindent
表現の変形理論はきわめて興味深いものとなりうる。

\begin{example}[群の表現]
\label{examples-defos-example-representations}
$\Gamma$ を群とする。圏 $\mathcal{F}$ を次のように定義する。
\begin{enumerate}
\item 対象は、$\mathcal{C}_\Lambda$ の対象 $A$、有限射影
$A$-加群 $M$、および準同型 $\rho : \Gamma \to \text{GL}_A(M)$
からなる三つ組 $(A, M, \rho)$ である。
\item 射 $(f, g) : (B, N, \tau) \to (A, M, \rho)$ は、
$\mathcal{C}_\Lambda$ における射 $f : B \to A$ と、$f$-線形かつ
$\Gamma$-同変であり、同型 $N \otimes_{B, f} A \cong M$ を誘導する
写像 $g : N \to M$ からなる。
\end{enumerate}
関手 $p : \mathcal{F} \to \mathcal{C}_\Lambda$ は $(A, M, \rho)$ を
$A$ に、$(f, g)$ を $f$ に送る。$p$ がグルーポイドをファイバーとする
余ファイバー圏であることは明らかである。有限次元 $k$-ベクトル空間
$V$ と表現 $\rho_0 : \Gamma \to \text{GL}_k(V)$ に対し、
$x_0 = (k, V, \rho_0)$ を対応する $\mathcal{F}(k)$ の対象とする。
次のようにおく。
$$
\Deformationcategory_{V, \rho_0} = \mathcal{F}_{x_0}
$$
\end{example}

\noindent
局所環上の有限射影加群はすべて有限自由であるから
（『代数』の補題 \ref{algebra-lemma-finite-projective}）、
次を得る。
$$
\begin{matrix}
\mathcal{F}(A)\text{ の対象の} \\
\text{同型類}
\end{matrix}
=
\coprod\nolimits_{n \geq 0}\quad
\begin{matrix}
\text{準同型 }\rho : \Gamma \to \text{GL}_n(A)\text{ の}\\
\text{GL}_n(A)\text{-共役類}
\end{matrix}
$$
これはすでに、節 \ref{examples-defos-section-finite-projective-modules} の議論よりも
興味深い。

\begin{lemma}
\label{examples-defos-lemma-representations-RS}
例 \ref{examples-defos-example-representations} は Rim--Schlessinger 条件 (RS) を満たす。
特に、任意の有限次元表現 $\rho_0 : \Gamma \to \text{GL}_k(V)$ に対して
$\Deformationcategory_{V, \rho_0}$ は変形圏である。
\end{lemma}

\begin{proof}
$A_1 \to A$ と $A_2 \to A$ を $\mathcal{C}_\Lambda$ の射とし、
$A_2 \to A$ は全射であると仮定する。『形式的変形理論』の補題
\ref{formal-defos-lemma-RS-2-categorical} により、関手
$\mathcal{F}(A_1 \times_A A_2) \to
\mathcal{F}(A_1) \times_{\mathcal{F}(A)} \mathcal{F}(A_2)$
が圏同値であることを示せば十分である。

\medskip\noindent
圏 $\mathcal{F}(A_1) \times_{\mathcal{F}(A)} \mathcal{F}(A_2)$ の対象
$$
((A_1, M_1, \rho_1), (A_2, M_2, \rho_2), (\text{id}_A, \varphi))
$$
を考える。補題 \ref{examples-defos-lemma-finite-projective-modules-RS} の証明で見たように、
有限射影 $A_1 \times_A A_2$-加群 $M_1 \times_\varphi M_2$ を考えられる。
$\varphi$ は与えられた作用と両立するので、次を得る。
$$
\rho_1 \times \rho_2 : \Gamma \longrightarrow
\text{GL}_{A_1 \times_A A_2}(M_1 \times_\varphi M_2)
$$
したがって $(M_1 \times_\varphi M_2, \rho_1 \times \rho_2)$ は
$\mathcal{F}(A_1 \times_A A_2)$ の対象である。この構成は上の関手の
擬逆関手を定める。
\end{proof}

\begin{lemma}
\label{examples-defos-lemma-representations-TI}
例 \ref{examples-defos-example-representations} において、
$\rho_0 : \Gamma \to \text{GL}_k(V)$ を有限次元表現とする。このとき
$$
T\Deformationcategory_{V, \rho_0} = \Ext^1_{k[\Gamma]}(V, V) =
H^1(\Gamma, \text{End}_k(V))
\;\text{and}\;
\text{Inf}(\Deformationcategory_{V, \rho_0}) = H^0(\Gamma, \text{End}_k(V))
$$
したがって $\text{Inf}(\Deformationcategory_{V, \rho_0})$ は常に有限次元であり、
$\Gamma$ が有限生成ならば $T\Deformationcategory_{V, \rho_0}$ も
有限次元である。
\end{lemma}

\begin{proof}
まず無限小自己同型を扱う。$M = V \otimes_k k[\epsilon]$ とし、
誘導される作用を $\rho_0' : \Gamma \to \text{GL}_n(M)$ とする。
補題 \ref{examples-defos-lemma-finite-projective-modules-TI} の証明と同様に、
無限小自己同型、すなわち
$\text{Inf}(\Deformationcategory_{V, \rho_0})$ の元は自己同型
$\gamma = \text{id} + \epsilon \psi : M \to M$ によって与えられる。
さらに $\psi$ は、$\rho_0$ による $\Gamma$ の作用と可換でなければならない。
したがって
$$
\text{Inf}(\Deformationcategory_{V, \rho_0}) = H^0(\Gamma, \text{End}_k(V))
$$
を得る。これは補題の主張どおりである。

\medskip\noindent
次に、$(k[\epsilon], M, \rho)$ を $k[\epsilon]$ 上の
$\mathcal{F}$ の対象とし、$\alpha : M \to V$ を、同型
$M/\epsilon M \to V$ を誘導する $\Gamma$-同変写像とする。
$M$ は $k[\epsilon]$-加群として自由なので、$\Gamma$-加群の拡大
$$
0 \to V \to M \xrightarrow{\alpha} V \to 0
$$
を得る。左側の写像の詳しい構成は省略する。逆に、上のような
$\Gamma$-加群の拡大が与えられれば、これを用いて $M$ に
$k[\epsilon]$-加群構造を入れ、上のような写像 $\alpha$ とともに
$\mathcal{F}(k[\epsilon])$ の対象を得る。したがって
$$
T\Deformationcategory_{V, \rho_0} = \Ext^1_{k[\Gamma]}(V, V)
$$
を得る。これは補題の主張どおりである。また『エタール・コホモロジー』の補題
\ref{etale-cohomology-lemma-ext-modules-hom} により、これは
$H^1(\Gamma, \text{End}_k(V))$ に等しい。

\medskip\noindent
次元に関する主張は『エタール・コホモロジー』の補題
\ref{etale-cohomology-lemma-finite-dim-group-cohomology} から従う。
\end{proof}

\noindent
例 \ref{examples-defos-example-representations} において、$\Gamma$ が有限生成であり、
$(V, \rho_0)$ が $k$ 上の $\Gamma$ の有限次元表現ならば、
$\Deformationcategory_{V, \rho_0}$ は $\mathcal{C}_\Lambda$ 上の関手からなる
滑らかなプロ表現可能グルーポイドによる表示をもち、したがってなおさら
（最小）半普遍的形式対象をもつ。これは補題
\ref{examples-defos-lemma-representations-RS}、\ref{examples-defos-lemma-representations-TI}、および
節 \ref{examples-defos-section-general} の一般的議論から従う。

\begin{lemma}
\label{examples-defos-lemma-representations-hull}
例 \ref{examples-defos-example-representations} において $\Gamma$ は有限生成と仮定する。
$\rho_0 : \Gamma \to \text{GL}_k(V)$ を有限次元表現とする。
$\Lambda$ は剰余体 $k$ をもつ完備局所環であると仮定する
（古典的な場合）。このとき関手
$$
F : \mathcal{C}_\Lambda \longrightarrow \textit{Sets},\quad
A \longmapsto \Ob(\Deformationcategory_{V, \rho_0}(A))/\cong
$$
は対象の同型類を対応させる関手であり、射影極限的包をもつ。

$H^0(\Gamma, \text{End}_k(V)) = k$ ならば、$F$ はプロ表現可能である。
\end{lemma}

\begin{proof}
射影極限的包の存在は、補題 \ref{examples-defos-lemma-representations-RS}、
\ref{examples-defos-lemma-representations-TI}、『形式的変形理論』の補題
\ref{formal-defos-lemma-RS-implies-S1-S2}、および注意
\ref{formal-defos-remark-compose-minimal-into-iso-classes} から従う。

\medskip\noindent
$H^0(\Gamma, \text{End}_k(V)) = k$ と仮定する。$F$ がプロ表現可能であることを
示すには、$F$ が変形関手であることを示せば十分である。
『形式的変形理論』の定理
\ref{formal-defos-theorem-Schlessinger-prorepresentability} を参照されたい。
言い換えれば、$F$ が (RS) を満たすことを示す必要がある。そのために
『形式的変形理論』の補題 \ref{formal-defos-lemma-RS-associated-functor}
の判定法を用いることができる。必要な自己同型群の全射性は、
$$
A \cdot \text{id}_M =
\text{End}_{A[\Gamma]}(M)
$$
が、$M \otimes_A k$ が $\Gamma$ の表現として $V$ と同型であるような
$\mathcal{F}$ の任意の対象 $(A, M, \rho)$ に対して成り立つことを示せば従う。
左辺は右辺に含まれるので、
$\text{length}_A \text{End}_{A[\Gamma]}(M) \leq \text{length}_A A$
を示せば十分である。互いに異なるイデアル
$(0) = I_n \subset \ldots \subset I_1 \subset A$
を $n = \text{length}(A)$ となるように選ぶ。これに応じて $M$ を
フィルター付けすると、$\Hom_{A[\Gamma]}(M, I_tM/I_{t + 1}M)$ の長さが
$1$ であることを示せば十分だと分かる。
$I_tM/I_{t + 1}M \cong M \otimes_A k$ であり、任意の
$A[\Gamma]$-加群写像 $M \to M \otimes_A k$ は商写像
$M \to M \otimes_A k$ を通じて一意に分解し、
$$
\text{End}_{A[\Gamma]}(M \otimes_A k) = \text{End}_{k[\Gamma]}(V) = k
$$
の元を与えるから、結論が従う。
\end{proof}



\section{連続表現}
\label{examples-defos-section-continuous-representations}

\noindent
無限 Galois 群をとり、その表現の変形理論を研究することは非常に興味深い。
\cite{Mazur-deforming} を参照されたい。

\begin{example}[位相群の表現]
\label{examples-defos-example-continuous-representations}
$\Gamma$ を位相群とする。圏 $\mathcal{F}$ を次のように定義する。
\begin{enumerate}
\item 対象は、$\mathcal{C}_\Lambda$ の対象 $A$、有限射影 $A$-加群 $M$、
および連続準同型 $\rho : \Gamma \to \text{GL}_A(M)$ からなる三つ組
$(A, M, \rho)$ である。ここで $\text{GL}_A(M)$ には離散位相を入れる。% P11-E0045
\footnote{別の定義として、$\rho$ が与える $G$-作用をもつ $A$-加群 $M$ が、
『エタール・コホモロジー』の定義
\ref{etale-cohomology-definition-G-module-continuous} にいう
$A\text{-}G$-加群であることを要求してもよい。しかし $M$ は有限
$A$-加群なので、これは同値である。}
\item 射 $(f, g) : (B, N, \tau) \to (A, M, \rho)$ は、
$\mathcal{C}_\Lambda$ における射 $f : B \to A$ と、$f$-線形かつ
$\Gamma$-同変であり、同型 $N \otimes_{B, f} A \cong M$ を誘導する
写像 $g : N \to M$ からなる。
\end{enumerate}
関手 $p : \mathcal{F} \to \mathcal{C}_\Lambda$ は $(A, M, \rho)$ を
$A$ に、$(f, g)$ を $f$ に送る。$p$ がグルーポイドをファイバーとする
余ファイバー圏であることは明らかである。有限次元 $k$-ベクトル空間
$V$ と連続表現 $\rho_0 : \Gamma \to \text{GL}_k(V)$ に対し、
$x_0 = (k, V, \rho_0)$ を対応する $\mathcal{F}(k)$ の対象とする。
次のようにおく。
$$
\Deformationcategory_{V, \rho_0} = \mathcal{F}_{x_0}
$$
\end{example}

\noindent
局所環上の有限射影加群はすべて有限自由であるから
（『代数』の補題 \ref{algebra-lemma-finite-projective}）、
次を得る。
$$
\begin{matrix}
\mathcal{F}(A)\text{ の対象の} \\
\text{同型類}
\end{matrix}
=
\coprod\nolimits_{n \geq 0}\quad
\begin{matrix}
\text{連続準同型 }\rho : \Gamma \to \text{GL}_n(A)\text{ の}\\
\text{GL}_n(A)\text{-共役類}
\end{matrix}
$$

\begin{lemma}
\label{examples-defos-lemma-continuous-representations-RS}
例 \ref{examples-defos-example-continuous-representations} は
Rim--Schlessinger 条件 (RS) を満たす。特に、任意の有限次元連続表現
$\rho_0 : \Gamma \to \text{GL}_k(V)$ に対して
$\Deformationcategory_{V, \rho_0}$ は変形圏である。
\end{lemma}

\begin{proof}
証明は補題 \ref{examples-defos-lemma-representations-RS} の証明とまったく同じである。
\end{proof}

\begin{lemma}
\label{examples-defos-lemma-continuous-representations-TI}
例 \ref{examples-defos-example-continuous-representations} において、
$\rho_0 : \Gamma \to \text{GL}_k(V)$ を有限次元連続表現とする。このとき
$$
T\Deformationcategory_{V, \rho_0} = H^1(\Gamma, \text{End}_k(V))
\quad\text{and}\quad
\text{Inf}(\Deformationcategory_{V, \rho_0}) = H^0(\Gamma, \text{End}_k(V))
$$
したがって $\text{Inf}(\Deformationcategory_{V, \rho_0})$ は常に有限次元であり、
$\Gamma$ が位相的有限生成ならば $T\Deformationcategory_{V, \rho_0}$ も
有限次元である。
\end{lemma}

\begin{proof}
証明は補題 \ref{examples-defos-lemma-representations-TI} の証明とまったく同じである。
\end{proof}

\noindent
例 \ref{examples-defos-example-continuous-representations} において、$\Gamma$ が位相的有限生成であり、
$(V, \rho_0)$ が $k$ 上の $\Gamma$ の有限次元連続表現ならば、
$\Deformationcategory_{V, \rho_0}$ は $\mathcal{C}_\Lambda$ 上の関手からなる
滑らかなプロ表現可能グルーポイドによる表示をもち、したがってなおさら
（最小）半普遍的形式対象をもつ。これは補題
\ref{examples-defos-lemma-continuous-representations-RS}、
\ref{examples-defos-lemma-continuous-representations-TI}、および節
\ref{examples-defos-section-general} の一般的議論から従う。

\begin{lemma}
\label{examples-defos-lemma-continuous-representations-hull}
例 \ref{examples-defos-example-continuous-representations} において、$\Gamma$ は位相的有限生成と
仮定する。$\rho_0 : \Gamma \to \text{GL}_k(V)$ を有限次元表現とする。
$\Lambda$ は剰余体 $k$ をもつ完備局所環であると仮定する
（古典的な場合）。このとき関手
$$
F : \mathcal{C}_\Lambda \longrightarrow \textit{Sets},\quad
A \longmapsto \Ob(\Deformationcategory_{V, \rho_0}(A))/\cong
$$
は対象の同型類を対応させる関手であり、射影極限的包をもつ。

$H^0(\Gamma, \text{End}_k(V)) = k$ ならば、$F$ はプロ表現可能である。
\end{lemma}

\begin{proof}
証明は補題 \ref{examples-defos-lemma-representations-hull} の証明とまったく同じである。
\end{proof}



\section{次数付き代数}
\label{examples-defos-section-graded-algebras}

\noindent
本節の例は、偏極付き固有スキームのスタックが代数スタックであることの
証明に用いる。このため、各斉次成分が有限射影加群である可換次数付き代数を
考える（この性質を「局所有限」と呼ぶこともある）。

\begin{example}[次数付き代数]
\label{examples-defos-example-graded-algebras}
圏 $\mathcal{F}$ を次のように定義する。
\begin{enumerate}
\item 対象は、$\mathcal{C}_\Lambda$ の対象 $A$ と、すべての $d \geq 0$
に対して $P_d$ が有限射影 $A$-加群であるような次数付き $A$-代数 $P$
からなる対 $(A, P)$ である。
\item 射 $(f, g) : (B, Q) \to (A, P)$ は、$\mathcal{C}_\Lambda$ における
射 $f : B \to A$ と、$f$-線形であり、同型
$Q \otimes_{B, f} A \cong P$ を誘導する写像 $g : Q \to P$ からなる。
\end{enumerate}
関手 $p : \mathcal{F} \to \mathcal{C}_\Lambda$ は $(A, P)$ を $A$ に、
$(f, g)$ を $f$ に送る。$p$ がグルーポイドをファイバーとする
余ファイバー圏であることは明らかである。すべての $d \geq 0$ に対して
$\dim_k(P_d) < \infty$ を満たす次数付き $k$-代数 $P$ に対し、
$x_0 = (k, P)$ を対応する $\mathcal{F}(k)$ の対象とする。次のようにおく。
$$
\Deformationcategory_P = \mathcal{F}_{x_0}
$$
\end{example}

\begin{lemma}
\label{examples-defos-lemma-graded-algebras-RS}
例 \ref{examples-defos-example-graded-algebras} は Rim--Schlessinger 条件 (RS) を満たす。
特に、任意の次数付き $k$-代数 $P$ に対して
$\Deformationcategory_P$ は変形圏である。
\end{lemma}

\begin{proof}
$A_1 \to A$ と $A_2 \to A$ を $\mathcal{C}_\Lambda$ の射とし、
$A_2 \to A$ は全射であると仮定する。『形式的変形理論』の補題
\ref{formal-defos-lemma-RS-2-categorical} により、関手
$\mathcal{F}(A_1 \times_A A_2) \to
\mathcal{F}(A_1) \times_{\mathcal{F}(A)} \mathcal{F}(A_2)$
が圏同値であることを示せば十分である。

\medskip\noindent
圏 $\mathcal{F}(A_1) \times_{\mathcal{F}(A)} \mathcal{F}(A_2)$ の対象
$$
((A_1, P_1), (A_2, P_2), (\text{id}_A, \varphi))
$$
を考える。$P_1 \times_\varphi P_2$ を考えよう。
$\varphi : P_1 \otimes_{A_1} A \to P_2 \otimes_{A_2} A$
は次数付き代数の同型なので、$P_1 \times_\varphi P_2$ の各次数成分は
有限射影 $A_1 \times_A A_2$-加群である。補題
\ref{examples-defos-lemma-finite-projective-modules-RS} の証明を参照されたい。
したがって $P_1 \times_\varphi P_2$ は
$\mathcal{F}(A_1 \times_A A_2)$ の対象である。この構成は上の関手の
擬逆関手を定め、証明は完了する。
\end{proof}

\begin{lemma}
\label{examples-defos-lemma-graded-algebras-TI}
例 \ref{examples-defos-example-graded-algebras} において、$P$ を次数付き $k$-代数とする。
このとき
$$
T\Deformationcategory_P
\quad\text{and}\quad
\text{Inf}(\Deformationcategory_P) = \text{Der}_k(P, P)
$$
$P$ が $k$ 上有限生成ならば、いずれも有限次元である。
\end{lemma}

\begin{proof}
まず無限小自己同型を扱う。$Q = P \otimes_k k[\epsilon]$ とおく。
$\text{Inf}(\Deformationcategory_P)$ の元は、上と同様に自己同型
$\gamma = \text{id} + \epsilon \delta : Q \to Q$ によって与えられる。
ここで $\delta : P \to P$ である。$\gamma$ が次数付きであることから
$\delta$ は次数 $0$ の斉次写像であり、$\gamma$ が $k$-線形であることから
$\delta$ は $k$-線形であり、$\gamma$ が乗法的であることから
$\delta$ は $k$-導分である。逆に、次数 $0$ の斉次 $k$-導分
$\delta : P \to P$ が与えられれば、上のような自己同型
$\gamma = \text{id} + \epsilon \delta$ を得る。したがって
$$
\text{Inf}(\Deformationcategory_P) = \text{Der}_k(P, P)
$$
を得る。これは補題の主張どおりである。明らかに、$P$ が次数成分 $P_i$
（$0 \leq i \leq N$）により生成されるならば、$\delta$ は
$0 \leq i \leq N$ に対する線形写像 $\delta_i : P_i \to P_i$ によって
決定される。したがって
$$
\dim_k \text{Der}_k(P, P) < \infty
$$
となり、所望の結論を得る。

\medskip\noindent
補題の証明を終えるため、変形空間が有限次元であることを示す。
そのため、次数付き $k$-代数の表示
$$
k[X_1, \ldots, X_n]/(F_1, \ldots, F_m) \longrightarrow P
$$
を選ぶ。ここで $\deg(X_i) = d_i$ であり、$F_j$ は次数 $e_j$ の斉次元である。
$Q$ を、各次数で有限自由な任意の次数付き $k[\epsilon]$-代数とし、
$(Q, \alpha)$ が $T\Deformationcategory_P$ の元を定めるような同型
$\alpha : Q/\epsilon Q \to P$ が与えられているとする。$P$ における
$X_i$ の像へ写る、次数 $d_i$ の斉次元 $q_i \in Q$ を選ぶ。このとき
$$
k[\epsilon][X_1, \ldots, X_n] \longrightarrow Q,\quad
X_i \longmapsto q_i
$$
を得る。$P = Q/\epsilon Q$ なので、Nakayama の補題によりこの写像は全射である。
簡単な図式追跡により、$Q$ では零に写り、$k[X_1, \ldots, X_n]$ では
$F_j$ に写る次数 $e_j$ の斉次元
$F_{\epsilon, j} \in k[\epsilon][X_1, \ldots, X_n]$ を選べる。このとき
$$
k[\epsilon][X_1, \ldots, X_n]/(F_{\epsilon, 1}, \ldots, F_{\epsilon, m})
\longrightarrow Q
$$
は、$Q$ の $k[\epsilon]$ 上の平坦性により $Q$ の表示である。次のように書く。
$$
F_{\epsilon, j} =  F_j + \epsilon G_j
$$
ベクトル $(G_1, \ldots, G_m)$ には若干の不定性がある。まず、
$F_{\epsilon, j}$ の選び方を変えることにより、$G_j$ を
$k[X_1, \ldots, X_n] \to P$ の核に属する次数 $e_j$ の任意の元だけ
変更できる。したがって $(G_1, \ldots, G_m)$ そのものではなく、元
$$
(g_1, \ldots, g_m) \in P_{e_1} \oplus \ldots \oplus P_{e_m}
$$
を記憶する。ここで $g_j$ は $P_{e_j}$ における $G_j$ の像である。
さらに $q_i$ の選択を、次数 $d_i$ の $p_i$ を用いて
$q_i + \epsilon p_i$ に変えると、計算（省略する）により $g_j$ は
$$
g_j^{new} = g_j - \sum\nolimits_{i = 1}^n p_i \partial F_j / \partial X_i
$$
に変わることが分かる。したがって $Q$ の同型類は、$k$-ベクトル空間
$$
W  = \Coker(P_{d_1} \oplus \ldots \oplus P_{d_n}
\xrightarrow{(\frac{\partial F_j}{\partial X_i})}
P_{e_1} \oplus \ldots \oplus P_{e_m})
$$
におけるベクトル $(G_1, \ldots, G_m)$ の像によって決定される。
このようにして単射
$$
T\Deformationcategory_P \longrightarrow W
$$
を得る。$W$ が明らかに有限次元であることから、補題の主張が従う。
\end{proof}

\noindent
例 \ref{examples-defos-example-graded-algebras} において $P$ が有限生成次数付き
$k$-代数ならば、$\Deformationcategory_P$ は $\mathcal{C}_\Lambda$ 上の
関手からなる滑らかなプロ表現可能グルーポイドによる表示をもち、
したがってなおさら（最小）半普遍的形式対象をもつ。これは補題
\ref{examples-defos-lemma-graded-algebras-RS}、\ref{examples-defos-lemma-graded-algebras-TI}、および
節 \ref{examples-defos-section-general} の一般的議論から従う。

\begin{lemma}
\label{examples-defos-lemma-graded-algebras-hull}
例 \ref{examples-defos-example-graded-algebras} において $P$ は有限生成次数付き
$k$-代数であると仮定する。また $\Lambda$ は剰余体 $k$ をもつ
完備局所環であると仮定する（古典的な場合）。このとき関手
$$
F : \mathcal{C}_\Lambda \longrightarrow \textit{Sets},\quad
A \longmapsto \Ob(\Deformationcategory_P(A))/\cong
$$
は対象の同型類を対応させる関手であり、射影極限的包をもつ。
\end{lemma}

\begin{proof}
これは補題 \ref{examples-defos-lemma-graded-algebras-RS}、
\ref{examples-defos-lemma-graded-algebras-TI}、『形式的変形理論』の補題
\ref{formal-defos-lemma-RS-implies-S1-S2}、および注意
\ref{formal-defos-remark-compose-minimal-into-iso-classes} から直ちに従う。
\end{proof}







\section{環}
\label{examples-defos-section-rings}

\noindent
環の変形理論はアフィンスキームの変形理論と同じである。環やスキームの
変形というときには、{\it 平坦}変形を考えているものとする。

\begin{example}[環]
\label{examples-defos-example-rings}
圏 $\mathcal{F}$ を次のように定義する。
\begin{enumerate}
\item 対象は、$\mathcal{C}_\Lambda$ の対象 $A$ と平坦 $A$-代数 $P$
からなる対 $(A, P)$ である。
\item 射 $(f, g) : (B, Q) \to (A, P)$ は、$\mathcal{C}_\Lambda$ における
射 $f : B \to A$ と、$f$-線形であり、同型
$Q \otimes_{B, f} A \cong P$ を誘導する写像 $g : Q \to P$ からなる。
\end{enumerate}
関手 $p : \mathcal{F} \to \mathcal{C}_\Lambda$ は $(A, P)$ を $A$ に、
$(f, g)$ を $f$ に送る。$p$ がグルーポイドをファイバーとする
余ファイバー圏であることは明らかである。$k$-代数 $P$ に対し、
$x_0 = (k, P)$ を対応する $\mathcal{F}(k)$ の対象とする。次のようにおく。
$$
\Deformationcategory_P = \mathcal{F}_{x_0}
$$
\end{example}

\begin{lemma}
\label{examples-defos-lemma-rings-RS}
例 \ref{examples-defos-example-rings} は Rim--Schlessinger 条件 (RS) を満たす。
特に、任意の $k$-代数 $P$ に対して $\Deformationcategory_P$ は変形圏である。
\end{lemma}

\begin{proof}
$A_1 \to A$ と $A_2 \to A$ を $\mathcal{C}_\Lambda$ の射とし、
$A_2 \to A$ は全射であると仮定する。『形式的変形理論』の補題
\ref{formal-defos-lemma-RS-2-categorical} により、関手
$\mathcal{F}(A_1 \times_A A_2) \to
\mathcal{F}(A_1) \times_{\mathcal{F}(A)} \mathcal{F}(A_2)$
が圏同値であることを示せば十分である。これは『代数詳論』の補題
\ref{more-algebra-lemma-properties-algebras-over-fibre-product}
の特別な場合である。
\end{proof}

\begin{lemma}
\label{examples-defos-lemma-rings-TI}
例 \ref{examples-defos-example-rings} において、$P$ を $k$-代数とする。このとき
$$
T\Deformationcategory_P = \text{Ext}^1_P(\NL_{P/k}, P)
\quad\text{and}\quad
\text{Inf}(\Deformationcategory_P) = \text{Der}_k(P, P)
$$
\end{lemma}

\begin{proof}
$\text{Inf}(\Deformationcategory_P)$ は、$P$ の $k[\epsilon]$ への自明変形
$P[\epsilon] = P \otimes_k k[\epsilon]$ の、$\epsilon$ を法として恒等写像に
等しい自己同型の集合であった。『変形理論』の補題
\ref{defos-lemma-huge-diagram} により、これは
$\Hom_P(\Omega_{P/k}, P)$ に等しく、さらに『代数』の補題
\ref{algebra-lemma-universal-omega} により $\text{Der}_k(P, P)$ に等しい。

\medskip\noindent
$T\Deformationcategory_P$ は、$P$ の $k[\epsilon]$ への平坦変形 $Q$ の
同型類の集合、より正確には $\Deformationcategory_P(k[\epsilon])$ の対象の
同型類の集合であった。また、$Q/\epsilon Q = P$ を満たす
$k[\epsilon]$-代数 $Q$ が $k[\epsilon]$ 上平坦であるための必要十分条件は、
$$
0 \to P \xrightarrow{\epsilon} Q \to P \to 0
$$
が完全であることである。これは『射の詳論』の補題
\ref{more-morphisms-lemma-deform}、より一般には『変形理論』の補題
\ref{defos-lemma-deform-module} で証明されている。したがって
『変形理論』の補題 \ref{defos-lemma-choices} を適用すると、
このような変形の同型類の集合が $\text{Ext}^1_P(\NL_{P/k}, P)$ に
等しいことが分かる。
\end{proof}

\begin{lemma}
\label{examples-defos-lemma-smooth}
例 \ref{examples-defos-example-rings} において、$P$ を滑らかな $k$-代数とする。このとき
$T\Deformationcategory_P = (0)$ である。
\end{lemma}

\begin{proof}
補題 \ref{examples-defos-lemma-rings-TI} により、
$\text{Ext}^1_P(\NL_{P/k}, P) = (0)$ を示せばよい。
$k \to P$ は滑らかなので、$\NL_{P/k}$ は、次数 $0$ に置かれた
有限射影 $P$-加群からなる複体と擬同型である。
\end{proof}

\begin{lemma}
\label{examples-defos-lemma-finite-type-rings-TI}
補題 \ref{examples-defos-lemma-rings-TI} において $P$ が有限型 $k$-代数ならば、
\begin{enumerate}
\item $\text{Inf}(\Deformationcategory_P)$ が有限次元であるための必要十分条件は
$\dim(P) = 0$ であることであり、
\item $\Spec(P) \to \Spec(k)$ が有限個の点を除いて滑らかならば、
$T\Deformationcategory_P$ は有限次元である。
\end{enumerate}
\end{lemma}

\begin{proof}
(1) の証明。$\text{Der}_k(P, P)$ を $P$-加群とみなす。これが $k$ 上
有限次元ならば、$P$-加群として有限長であり、したがってその台は
$\Spec(P)$ の有限個の閉点に含まれる
（『代数』の補題 \ref{algebra-lemma-simple-pieces}）。
$\text{Der}_k(P, P) = \Hom_P(\Omega_{P/k}, P)$ なので、任意の素イデアル
$\mathfrak p \subset P$ に対して
$\text{Der}_k(P, P)_\mathfrak p = \text{Der}_k(P_\mathfrak p, P_\mathfrak p)$
を得る（ここでは『代数』の補題
\ref{algebra-lemma-differentials-localize}、
\ref{algebra-lemma-differentials-finitely-presented}、および
\ref{algebra-lemma-hom-from-finitely-presented} を用いる）。
$\mathfrak p$ を、次元 $d > 0$ の既約成分に対応する $P$ の極小素イデアルとする。
このとき $P_\mathfrak p$ は $k$ 上本質的に有限型で剰余体をもつ
Artin 局所環であり、% P11-E0047
たとえば『代数』の補題
\ref{algebra-lemma-characterize-smooth-over-field} により
$\Omega_{P_\mathfrak p/k}$ は零でない。Artin 局所環上の任意の非零有限加群は、
剰余体と同型な部分加群と商加群をともにもつ。したがって
$\text{Der}_k(P_\mathfrak p, P_\mathfrak p) =
\Hom_{P_\mathfrak p}(\Omega_{P_\mathfrak p/k}, P_\mathfrak p)$
も零でない。以上を合わせると (1) が従う。

\medskip\noindent
(2) の証明。$P$ の素イデアル $\mathfrak p$ に対して
$\NL_{P_\mathfrak p/k} = (\NL_{P/k})_\mathfrak p$
（『代数』の補題 \ref{algebra-lemma-localize-NL}）および
$\text{Ext}_P^1(\NL_{P/k}, P)_\mathfrak p =
\text{Ext}_{P_\mathfrak p}^1(\NL_{P_\mathfrak p/k}, P_\mathfrak p)$
（『代数詳論』の補題
\ref{more-algebra-lemma-pseudo-coherence-and-base-change-ext}）を用いる。
素イデアル $\mathfrak p \subset P$ に対し、$k \to P$ が $\mathfrak p$ で
滑らかであるための必要十分条件は、$(\NL_{P/k})_\mathfrak p$ が次数 $0$ に
置かれた有限射影加群と擬同型であることである（これは滑らかな環準同型の
定義から直ちに従うが、より強い『代数』の補題
\ref{algebra-lemma-smooth-at-point} からも従う）。

\medskip\noindent
$P$ は有限個の素イデアルを除いて $k$ 上滑らかであると仮定する。
この「悪い」素イデアルたちは、$\Spec(P)$ の閉部分集合をなすことと
『代数』の補題 \ref{algebra-lemma-finite-type-algebra-finite-nr-primes}
により、極大イデアル $\mathfrak m_1, \ldots, \mathfrak m_n \subset P$ である。
$\mathfrak p \not \in \{\mathfrak m_1, \ldots, \mathfrak m_n\}$ に対して、
上の結果から $\text{Ext}^1_P(\NL_{P/k}, P)_\mathfrak p = 0$ である。
したがって $\text{Ext}^1_P(\NL_{P/k}, P)$ は、その台が
$\{\mathfrak m_1, \ldots, \mathfrak m_r\}$ に含まれる有限 $P$-加群である。% P11-E0048
たとえば『代数』の命題
\ref{algebra-proposition-minimal-primes-associated-primes} により、
$\text{Ext}^1_P(\NL_{P/k}, P)$ の $k$ 上の次元は
$\dim_k \kappa(\mathfrak m_i)$ の有限な整係数一次結合である。したがって
Hilbert の零点定理（『代数』の定理
\ref{algebra-theorem-nullstellensatz}）により有限である。
\end{proof}

\noindent
例 \ref{examples-defos-example-rings} において、$P$ を有限型 $k$-代数とする。
$\Deformationcategory_P$ が $\mathcal{C}_\Lambda$ 上の関手からなる
滑らかなプロ表現可能グルーポイドによる表示をもつための必要十分条件は
$\dim(P) = 0$ である。さらに、$\Spec(P) \to \Spec(k)$ の特異点が
有限個ならば、$\Deformationcategory_P$ は半普遍的形式対象をもつ。
これは補題 \ref{examples-defos-lemma-rings-RS}、\ref{examples-defos-lemma-finite-type-rings-TI}、および
節 \ref{examples-defos-section-general} の一般的議論から従う。

\begin{lemma}
\label{examples-defos-lemma-rings-hull}
例 \ref{examples-defos-example-rings} において、$P$ は有限型 $k$-代数であり、
$\Spec(P) \to \Spec(k)$ は有限個の点を除いて滑らかであると仮定する。
また $\Lambda$ は剰余体 $k$ をもつ完備局所環であると仮定する
（古典的な場合）。このとき関手
$$
F : \mathcal{C}_\Lambda \longrightarrow \textit{Sets},\quad
A \longmapsto \Ob(\Deformationcategory_P(A))/\cong
$$
は対象の同型類を対応させる関手であり、射影極限的包をもつ。
\end{lemma}

\begin{proof}
これは補題 \ref{examples-defos-lemma-rings-RS}、\ref{examples-defos-lemma-finite-type-rings-TI}、
『形式的変形理論』の補題 \ref{formal-defos-lemma-RS-implies-S1-S2}、
および注意 \ref{formal-defos-remark-compose-minimal-into-iso-classes}
から直ちに従う。
\end{proof}

\begin{lemma}
\label{examples-defos-lemma-localization}
例 \ref{examples-defos-example-rings} において、$P$ を $k$-代数とし、
$S \subset P$ を乗法的部分集合とする。変形圏の自然な関手
$$
\Deformationcategory_P \longrightarrow \Deformationcategory_{S^{-1}P}
$$
が存在する。
\end{lemma}

\begin{proof}
$P$ の変形を局所化すれば、その局所化の変形が得られる。これは明らかなので、
読者には証明を飛ばすことを勧める。より正確には、
$(A, Q) \to (k, P)$ を $\mathcal{F}$ の射、すなわち
$\Deformationcategory_P$ の対象とする。$S_Q \subset Q$ を $S$ の逆像とする。% P11-E0049
したがって $(A, S_Q^{-1}Q) \to (k, S^{-1}P)$ は
$\Deformationcategory_{S^{-1}P}$ の所望の対象である。
\end{proof}

\begin{lemma}
\label{examples-defos-lemma-henselization}
例 \ref{examples-defos-example-rings} において、$P$ を $k$-代数とし、
$J \subset P$ をイデアルとする。$(P^h, J^h)$ を対 $(P, J)$ の
ヘンゼル化とする。変形圏の自然な関手
$$
\Deformationcategory_P \longrightarrow \Deformationcategory_{P^h}
$$
が存在する。
\end{lemma}

\begin{proof}
$P$ の変形をヘンゼル化すれば、そのヘンゼル化の変形が得られる。
これは明らかなので、読者には証明を飛ばすことを勧める。より正確には、
$(A, Q) \to (k, P)$ を $\mathcal{F}$ の射、すなわち
$\Deformationcategory_P$ の対象とする。$J_Q \subset Q$ を $Q$ における
$J$ の逆像とし、$(Q^h, J_Q^h)$ を対 $(Q, J_Q)$ のヘンゼル化とする。
$Q \to Q^h$ は平坦であるから（『代数詳論』の補題
\ref{more-algebra-lemma-henselization-flat}）、$Q^h$ は $A$ 上平坦である。
『代数詳論』の補題 \ref{more-algebra-lemma-henselization-integral} により、
写像 $Q^h \to P^h$ は同型
$Q^h \otimes_A k = Q^h \otimes_Q P = P^h$.
を誘導する。したがって $(A, Q^h) \to (k, P^h)$ は
$\Deformationcategory_{P^h}$ の所望の対象である。
\end{proof}

\begin{lemma}
\label{examples-defos-lemma-strict-henselization}
例 \ref{examples-defos-example-rings} において、$P$ を $k$-代数とする。
$P$ は局所環であると仮定し、$P^{sh}$ を $P$ の強ヘンゼル化とする。
変形圏の自然な関手
$$
\Deformationcategory_P \longrightarrow \Deformationcategory_{P^{sh}}
$$
が存在する。
\end{lemma}

\begin{proof}
$P$ の変形を強ヘンゼル化すれば、その強ヘンゼル化の変形が得られる。
これは明らかなので、読者には証明を飛ばすことを勧める。より正確には、
$(A, Q) \to (k, P)$ を $\mathcal{F}$ の射、すなわち
$\Deformationcategory_P$ の対象とする。全射 $Q \to P$ の核は冪零なので、
$Q$ は $P$ と同じ剰余体をもつ局所環である。$Q^{sh}$ を $Q$ の
強ヘンゼル化とする。$Q \to Q^{sh}$ は平坦であるから
（『代数詳論』の補題
\ref{more-algebra-lemma-dumb-properties-henselization}）、
$Q^{sh}$ は $A$ 上平坦である。『代数』の補題
\ref{algebra-lemma-quotient-strict-henselization} により、
写像 $Q^{sh} \to P^{sh}$ は同型
$Q^{sh} \otimes_A k = Q^{sh} \otimes_Q P = P^{sh}$.
を誘導する。したがって $(A, Q^{sh}) \to (k, P^{sh})$ は
$\Deformationcategory_{P^{sh}}$ の所望の対象である。
\end{proof}

\begin{lemma}
\label{examples-defos-lemma-completion}
例 \ref{examples-defos-example-rings} において、$P$ を $k$-代数とする。
$P$ はネーター的であると仮定し、$J \subset P$ をイデアルとする。% P11-E0050
$P^\wedge$ を $J$-進完備化とする。変形圏の自然な関手
$$
\Deformationcategory_P \longrightarrow \Deformationcategory_{P^\wedge}
$$
が存在する。
\end{lemma}

\begin{proof}
$P$ の変形を完備化すれば、その完備化の変形が得られる。これは明らかなので、
読者には証明を飛ばすことを勧める。より正確には、
$(A, Q) \to (k, P)$ を $\mathcal{F}$ の射、すなわち
$\Deformationcategory_P$ の対象とする。$Q$ はネーター環であることに注意する。
実際、全射環準同型 $Q \to P$ の核は冪零かつ有限生成であり、$P$ は
ネーター的であるから、『代数』の補題
\ref{algebra-lemma-completion-Noetherian} を適用できる。
$J_Q \subset Q$ を $Q$ における $J$ の逆像とし、$Q^\wedge$ を $Q$ の
$J_Q$-進完備化とする。$Q \to Q^\wedge$ は平坦であるから
（『代数』の補題 \ref{algebra-lemma-completion-flat}）、
$Q^\wedge$ は $A$ 上平坦である。誘導写像 $Q^\wedge \to P^\wedge$ は、
たとえば『代数』の補題 \ref{algebra-lemma-completion-tensor} により同型
$Q^\wedge \otimes_A k = Q^\wedge \otimes_Q P = P^\wedge$
を誘導する。したがって $(A, Q^\wedge) \to (k, P^\wedge)$ は
$\Deformationcategory_{P^\wedge}$ の所望の対象である。
\end{proof}

\begin{lemma}
\label{examples-defos-lemma-power-series-rings-TI}
補題 \ref{examples-defos-lemma-rings-TI} において、ある非零元
$f \in (x_1, \ldots, x_n)^2$ に対して
$P = k[[x_1, \ldots, x_n]]/(f)$ であるとする。このとき
\begin{enumerate}
\item $\text{Inf}(\Deformationcategory_P)$ が有限次元であるための
必要十分条件は $n = 1$ であり、
\item 次が成り立つならば $T\Deformationcategory_P$ は有限次元である。
$$
\sqrt{(f, \partial f/\partial x_1, \ldots,  \partial f/\partial x_n)} =
(x_1, \ldots, x_n)
$$
\end{enumerate}
\end{lemma}

\begin{proof}
(1) の証明。$k[[x_1, \ldots, x_n]]$ の $k$ 上の導分
$\partial/\partial x_i$ を考え、$f_i = \partial f/\partial x_i$ と書く。
$k[[x_1, \ldots, x_n]]$ の導分
$$
\theta = \sum h_i \partial/\partial x_i
$$
が $P = k[[x_1, \ldots, x_n]]/(f)$ の導分を誘導するための必要十分条件は
$\sum h_i f_i \in (f)$ である。さらに、誘導される $P$ の導分が零である
ための必要十分条件は、$i = 1, \ldots, n$ に対して $h_i \in (f)$
である。したがって
$$
\Ker((f_1, \ldots, f_n) : P^{\oplus n} \longrightarrow P) \subset
\text{Der}_k(P, P)
$$
を得る。左辺が有限次元 $k$-ベクトル空間となるのは $n = 1$ の場合に限る。
証明は省略する。また $n = 1$ ならば右辺が有限次元であることの確認も
読者に委ねる。これで (1) が証明された。

\medskip\noindent
(2) の証明。補題 \ref{examples-defos-lemma-rings-TI} の証明と同様に、$Q$ を
$P$ の $k[\epsilon]$ 上の平坦変形とする。$P$ における $x_i$ の像の
持ち上げ $q_i \in Q$ を選ぶ。このとき $Q$ は、極大イデアルが
$q_1, \ldots, q_n$ と $\epsilon$ で生成される完備局所環である
（簡単な議論は省略する）。したがって全射
$$
k[\epsilon][[x_1, \ldots, x_n]] \longrightarrow Q,\quad
x_i \longmapsto q_i
$$
を得る。$Q$ で零に写る形
$f + \epsilon g \in k[\epsilon][[x_1, \ldots, x_n]]$ の元を選ぶ。
$g$ は $(f)$ を法としてよく定まることに注意する。$Q$ は
$k[\epsilon]$ 上平坦なので
$$
Q = k[\epsilon][[x_1, \ldots, x_n]]/(f + \epsilon g)
$$
を得る。最後に、$q_i$ の選び方を変えることは、ある
$h_i \in k[[x_1, \ldots, x_n]]$ に対して座標 $x_i$ を
$x_i + \epsilon h_i$ に変えることに相当する。% P11-E0051
このとき $f_i = \partial f/\partial x_i$ とおけば、
$f + \epsilon g$ は $f + \epsilon (g + \sum h_i f_i)$ に変わる。
したがって変形 $Q$ の同型類は、次の商の元によって決定される。
$$
k[[x_1, \ldots, x_n]]/
(f, \partial f/\partial x_1, \ldots,  \partial f/\partial x_n)
$$
これが $k$ 上有限次元であるための必要十分条件は、その台が
$k[[x_1, \ldots, x_n]]$ の閉点であることであり、さらにこれは
$\sqrt{(f, \partial f/\partial x_1, \ldots,  \partial f/\partial x_n)} =
(x_1, \ldots, x_n)$.
\end{proof}






\section{スキーム}
\label{examples-defos-section-schemes}

\noindent
スキームの変形理論を扱う。

\begin{example}[スキーム]
\label{examples-defos-example-schemes}
圏 $\mathcal{F}$ を次のように定義する。
\begin{enumerate}
\item 対象は、$\mathcal{C}_\Lambda$ の対象 $A$ と $A$ 上平坦なスキーム $X$
からなる対 $(A, X)$ である。
\item 射 $(f, g) : (B, Y) \to (A, X)$ は、$\mathcal{C}_\Lambda$ における
射 $f : B \to A$ と、次がスキームのデカルト可換図式となるような射
$g : X \to Y$ からなる。
$$
\xymatrix{
X \ar[r]_g \ar[d] & Y \ar[d] \\
\Spec(A) \ar[r]^f & \Spec(B)
}
$$
\end{enumerate}
関手 $p : \mathcal{F} \to \mathcal{C}_\Lambda$ は $(A, X)$ を $A$ に、
$(f, g)$ を $f$ に送る。$p$ がグルーポイドをファイバーとする
余ファイバー圏であることは明らかである。$k$ 上のスキーム $X$ に対し、
$x_0 = (k, X)$ を対応する $\mathcal{F}(k)$ の対象とする。次のようにおく。
$$
\Deformationcategory_X = \mathcal{F}_{x_0}
$$
\end{example}

\begin{lemma}
\label{examples-defos-lemma-schemes-RS}
例 \ref{examples-defos-example-schemes} は Rim--Schlessinger 条件 (RS) を満たす。
特に、任意の $k$ 上のスキーム $X$ に対して
$\Deformationcategory_X$ は変形圏である。
\end{lemma}

\begin{proof}
$A_1 \to A$ と $A_2 \to A$ を $\mathcal{C}_\Lambda$ の射とし、
$A_2 \to A$ は全射であると仮定する。『形式的変形理論』の補題
\ref{formal-defos-lemma-RS-2-categorical} により、関手
$\mathcal{F}(A_1 \times_A A_2) \to
\mathcal{F}(A_1) \times_{\mathcal{F}(A)} \mathcal{F}(A_2)$
が圏同値であることを示せば十分である。図式
$$
\xymatrix{
\Spec(A) \ar[r] \ar[d] & \Spec(A_2) \ar[d] \\
\Spec(A_1) \ar[r] &
\Spec(A_1 \times_A A_2)
}
$$
は『射の詳論』の補題
\ref{more-morphisms-lemma-pushout-along-thickening} にいう押し出し図式である。
したがって本補題は『射の詳論』の補題
\ref{more-morphisms-lemma-equivalence-categories-schemes-over-pushout-flat}
の特別な場合である。
\end{proof}

\begin{lemma}
\label{examples-defos-lemma-schemes-TI}
例 \ref{examples-defos-example-schemes} において、$X$ を $k$ 上のスキームとする。このとき
$$
\text{Inf}(\Deformationcategory_X) =
\text{Ext}^0_{\mathcal{O}_X}(\NL_{X/k}, \mathcal{O}_X) =
\Hom_{\mathcal{O}_X}(\Omega_{X/k}, \mathcal{O}_X) =
\text{Der}_k(\mathcal{O}_X, \mathcal{O}_X)
$$
かつ
$$
T\Deformationcategory_X =
\text{Ext}^1_{\mathcal{O}_X}(\NL_{X/k}, \mathcal{O}_X)
$$
\end{lemma}

\begin{proof}
$\text{Inf}(\Deformationcategory_X)$ は、$X$ の $k[\epsilon]$ への自明変形
$X' = X \times_{\Spec(k)} \Spec(k[\epsilon])$ の自己同型で、
$\epsilon$ を法として恒等写像となるもの全体の集合であることを思い出そう。
『変形理論』の補題 \ref{defos-lemma-deform} により、これは
$\text{Ext}^0_{\mathcal{O}_X}(\NL_{X/k}, \mathcal{O}_X)$ に等しい。
等式 $\text{Ext}^0_{\mathcal{O}_X}(\NL_{X/k}, \mathcal{O}_X) =
\Hom_{\mathcal{O}_X}(\Omega_{X/k}, \mathcal{O}_X)$ は、
『射の詳論』の補題
\ref{more-morphisms-lemma-netherlander-quasi-coherent} から従う。
また、等式
$\Hom_{\mathcal{O}_X}(\Omega_{X/k}, \mathcal{O}_X) =
\text{Der}_k(\mathcal{O}_X, \mathcal{O}_X)$
は『射』の補題
\ref{morphisms-lemma-universal-derivation-universal} から従う。

\medskip\noindent
$T_{x_0}\Deformationcategory_X$ は、$X$ の $k[\epsilon]$ への平坦変形
$X'$ の同型類全体の集合、より正確には
$\Deformationcategory_X(k[\epsilon])$ の同型類全体の集合であることを思い出そう。
したがって補題の第2の主張は、『変形理論』の補題
\ref{defos-lemma-deform} から従う。
\end{proof}

\begin{lemma}
\label{examples-defos-lemma-proper-schemes-TI}
補題 \ref{examples-defos-lemma-schemes-TI} において $X$ が $k$ 上固有ならば、
$\text{Inf}(\Deformationcategory_X)$ と $T\Deformationcategory_X$ は
有限次元である。
\end{lemma}

\begin{proof}
同補題により、
$\Ext^1_{\mathcal{O}_X}(\NL_{X/k}, \mathcal{O}_X)$ と
$\Ext^0_{\mathcal{O}_X}(\NL_{X/k}, \mathcal{O}_X)$ が有限次元であることを
示せばよい。『射の詳論』の補題
\ref{more-morphisms-lemma-netherlander-fp}
と $X$ がネーターであることから、$\NL_{X/k}$ のコホモロジー層は
連接であり、次数 $0$ と $-1$ を除いて零であることが分かる。
『スキームの導来圏』の補題 \ref{perfect-lemma-ext-finite} により、
上に表示した $\Ext$ 群は有限次元 $k$ ベクトル空間である。これで証明が完了する。
\end{proof}

\noindent
例 \ref{examples-defos-example-schemes} において $X$ が $k$ 上の固有スキームならば、
$\Deformationcategory_X$ は $\mathcal{C}_\Lambda$ 上の関手のなす
滑らかなプロ表現可能グルーポイドによる表示をもち、したがって特に
（極小）半普遍的形式対象をもつ。これは補題 \ref{examples-defos-lemma-schemes-RS}、
\ref{examples-defos-lemma-proper-schemes-TI}
および節 \ref{examples-defos-section-general} の一般的な議論から従う。

\begin{lemma}
\label{examples-defos-lemma-schemes-hull}
例 \ref{examples-defos-example-schemes} において $X$ を固有 $k$ スキームとする。
$\Lambda$ を剰余体 $k$ をもつ完備局所環と仮定する（古典的な場合）。
このとき関手
$$
F : \mathcal{C}_\Lambda \longrightarrow \textit{Sets},\quad
A \longmapsto \Ob(\Deformationcategory_X(A))/\cong
$$
は、対象の同型類を対応させる関手であり、射影極限的包をもつ。さらに
$\text{Der}_k(\mathcal{O}_X, \mathcal{O}_X) = 0$ ならば、
$F$ はプロ表現可能である。
\end{lemma}

\begin{proof}
射影極限的包の存在は、補題 \ref{examples-defos-lemma-schemes-RS}、
\ref{examples-defos-lemma-proper-schemes-TI}、『形式的変形理論』の補題
\ref{formal-defos-lemma-RS-implies-S1-S2}、および注意
\ref{formal-defos-remark-compose-minimal-into-iso-classes} から直ちに従う。

\medskip\noindent
$\text{Der}_k(\mathcal{O}_X, \mathcal{O}_X) = 0$ と仮定する。
このとき『形式的変形理論』の補題
\ref{formal-defos-lemma-infdef-trivial} により、
$\Deformationcategory_X$ と $F$ は同値である。したがって $F$ は
有限次元接空間をもつ変形関手であり（$\Deformationcategory_X$ は変形圏である）、
『形式的変形理論』の定理
\ref{formal-defos-theorem-Schlessinger-prorepresentability} を適用できる。
\end{proof}

\begin{lemma}
\label{examples-defos-lemma-open}
例 \ref{examples-defos-example-schemes} において $X$ を $k$ 上のスキームとし、
$U \subset X$ を開部分スキームとする。変形圏の自然な関手
$$
\Deformationcategory_X \longrightarrow \Deformationcategory_U
$$
が存在する。
\end{lemma}

\begin{proof}
$X$ の変形が与えられたとき、その対応する開部分を取れば $U$ の変形を得る。
詳細は省略する。
\end{proof}

\begin{lemma}
\label{examples-defos-lemma-affine}
例 \ref{examples-defos-example-schemes} において $X = \Spec(P)$ を $k$ 上の
アフィンスキームとする。例 \ref{examples-defos-example-rings} の
$\Deformationcategory_P$ に対して、変形圏の自然な同値
$$
\Deformationcategory_X \longrightarrow \Deformationcategory_P
$$
が存在する。
\end{lemma}

\begin{proof}
この関手は $(A, Y)$ を $\Gamma(Y, \mathcal{O}_Y)$ に送る。
『射の詳論』の補題
\ref{more-morphisms-lemma-thickening-affine-scheme} により、
$X$ の任意の変形はアフィンなので、これは確かに定義できる。
\end{proof}

\begin{lemma}
\label{examples-defos-lemma-local-ring}
例 \ref{examples-defos-example-schemes} において $X$ を $k$ 上のスキームとし、
$p \in X$ を点とする。例 \ref{examples-defos-example-rings} の
$\Deformationcategory_{\mathcal{O}_{X, p}}$ に対して、変形圏の自然な関手
$$
\Deformationcategory_X
\longrightarrow
\Deformationcategory_{\mathcal{O}_{X, p}}
$$
が存在する。
\end{lemma}

\begin{proof}
$p$ を含むアフィン開部分 $U = \Spec(P) \subset X$ を選ぶ。
このとき $\mathcal{O}_{X, p}$ は $P$ の局所化である。補題
\ref{examples-defos-lemma-open}、\ref{examples-defos-lemma-affine}、\ref{examples-defos-lemma-localization}
の関手を合成すればよい。
\end{proof}

\begin{situation}
\label{examples-defos-situation-glueing}
$\Lambda \to k$ を節 \ref{examples-defos-section-general} のものとする。
$X$ を $k$ 上のスキームとし、アフィン開被覆
$X = U_1 \cup U_2$ をもち、$U_{12} = U_1 \cap U_2$ もアフィンであるとする。
$U_1 = \Spec(P_1)$、$U_2 = \Spec(P_2)$、
$U_{12} = \Spec(P_{12})$ と書く。
$\Deformationcategory_X$、$\Deformationcategory_{U_1}$、
$\Deformationcategory_{U_2}$、$\Deformationcategory_{U_{12}}$ は
例 \ref{examples-defos-example-schemes} のものとし、
$\Deformationcategory_{P_1}$、$\Deformationcategory_{P_2}$、
$\Deformationcategory_{P_{12}}$ は例 \ref{examples-defos-example-rings} のものとする。
\end{situation}

\begin{lemma}
\label{examples-defos-lemma-glueing}
状況 \ref{examples-defos-situation-glueing} において、変形圏の同値
$$
\Deformationcategory_X =
\Deformationcategory_{P_1}
\times_{\Deformationcategory_{P_{12}}}
\Deformationcategory_{P_2}
$$
が存在する。例 \ref{examples-defos-example-schemes} および \ref{examples-defos-example-rings} を参照せよ。
\end{lemma}

\begin{proof}
補題 \ref{examples-defos-lemma-open} の関手が同値
$$
\Deformationcategory_X \longrightarrow
\Deformationcategory_{U_1}
\times_{\Deformationcategory_{U_{12}}}
\Deformationcategory_{U_2}
$$
を定めることを示せば十分である。実際、そのとき補題
\ref{examples-defos-lemma-affine} を適用して環の言葉に移せる。これを示すため、
擬逆を構成する。補題 \ref{examples-defos-lemma-open} の関手を
$F_i : \Deformationcategory_{U_i} \to \Deformationcategory_{U_{12}}$
と書く。右辺の対象は、$\mathcal{C}_\Lambda$ の対象 $A$、対象
$(A, V_1) \to (k, U_1)$ と $(A, V_2) \to (k, U_2)$、および射
$$
g : F_1(A, V_1) \to F_2(A, V_2)
$$
によって与えられる。ここで $F_i(A, V_i) = (A, V_{i, 3 - i})$ であり、
$V_{i, 3 - i} \subset V_i$ は、$k$ への基底変換が
$U_{12} \subset U_i$ となる開部分スキームである。射 $g$ は、$k$ 上の
$\text{id} : U_{12} \to U_{12}$ と両立する $A$ 上のスキームの同型
$V_{1, 2} \to V_{2, 1}$ を定める。したがって
$(\{1, 2\}, V_i, V_{i, 3 - i}, g, g^{-1})$ は、『スキーム』の節
\ref{schemes-section-glueing-schemes} にいう貼り合わせデータである。
$Y$ をその貼り合わせとする（『スキーム』の補題
\ref{schemes-lemma-glue} を参照せよ）。すると $Y$ は $A$ 上のスキームであり、
上記の両立性から標準的な同型
$Y \times_{\Spec(A)} \Spec(k) = X$.
が得られる。したがって $(A, Y) \to (k, X)$ は
$\Deformationcategory_X$ の対象である。この構成が関手を定め、
与えられた関手の擬逆であることの確認は省略する。
\end{proof}






\section{スキームの射}
\label{examples-defos-section-schemes-morphisms}

\noindent
スキームの射の変形理論を扱う。
もちろん、これはスキームの図式の変形の一例にすぎない。

\begin{example}[スキームの射]
\label{examples-defos-example-schemes-morphisms}
圏 $\mathcal{F}$ を次のように定義する。
\begin{enumerate}
\item 対象は、$\mathcal{C}_\Lambda$ の対象 $A$ と $A$ 上のスキームの射
$X \to Y$ からなる対 $(A, X \to Y)$ であって、$X$ と $Y$ がともに
$A$ 上平坦なものである。
\item 射 $(f, g, h) : (A', X' \to Y') \to (A, X \to Y)$ は、
$\mathcal{C}_\Lambda$ における射 $f : A' \to A$ と、次が両方の正方形が
デカルトであるスキームの可換図式となるようなスキームの射
$g : X \to X'$ および $h : Y \to Y'$ からなる。
$$
\xymatrix{
X \ar[r]_g \ar[d] & X' \ar[d] \\
Y \ar[r]_h \ar[d] & Y' \ar[d] \\
\Spec(A) \ar[r]^f & \Spec(A')
}
$$
\end{enumerate}
関手 $p : \mathcal{F} \to \mathcal{C}_\Lambda$ は $(A, X \to Y)$ を
$A$ に、$(f, g, h)$ を $f$ に送る。$p$ がグルーポイドをファイバーとする
余ファイバー圏であることは明らかである。$k$ 上のスキームの射
$X \to Y$ に対し、$x_0 = (k, X \to Y)$ を対応する
$\mathcal{F}(k)$ の対象とする。次のようにおく。
$$
\Deformationcategory_{X \to Y} = \mathcal{F}_{x_0}
$$
\end{example}

\begin{lemma}
\label{examples-defos-lemma-schemes-morphisms-RS}
例 \ref{examples-defos-example-schemes-morphisms} は Rim--Schlessinger 条件 (RS) を満たす。
特に、$k$ 上の任意のスキームの射 $X \to Y$ に対して
$\Deformationcategory_{X \to Y}$ は変形圏である。
\end{lemma}

\begin{proof}
$A_1 \to A$ と $A_2 \to A$ を $\mathcal{C}_\Lambda$ の射とし、
$A_2 \to A$ は全射であると仮定する。『形式的変形理論』の補題
\ref{formal-defos-lemma-RS-2-categorical}
により、関手
$\mathcal{F}(A_1 \times_A A_2) \to
\mathcal{F}(A_1) \times_{\mathcal{F}(A)} \mathcal{F}(A_2)$
が圏同値であることを示せば十分である。図式
$$
\xymatrix{
\Spec(A) \ar[r] \ar[d] & \Spec(A_2) \ar[d] \\
\Spec(A_1) \ar[r] &
\Spec(A_1 \times_A A_2)
}
$$
は『射の詳論』の補題
\ref{more-morphisms-lemma-pushout-along-thickening} にいう押し出し図式である。
したがって本補題は『射の詳論』の補題
\ref{more-morphisms-lemma-equivalence-categories-schemes-over-pushout-flat}
から直ちに従う。実際、この補題は $A_1 \times_A A_2$ 上平坦なスキームの圏を、
$A$ 上平坦なスキームの圏の上での、$A_1$ 上平坦なスキームの圏と
$A_2$ 上平坦なスキームの圏とのファイバー積として記述している。
\end{proof}

\begin{lemma}
\label{examples-defos-lemma-schemes-morphisms-TI}
例 \ref{examples-defos-example-schemes} において、$f : X \to Y$ を $k$ 上のスキームの射とする。
% P11-E0052
$k$ ベクトル空間の次の標準的な完全列が存在する。
$$
\xymatrix{
0 \ar[r] &
\text{Inf}(\Deformationcategory_{X \to Y}) \ar[r] &
\text{Inf}(\Deformationcategory_X \times \Deformationcategory_Y) \ar[r] &
\text{Der}_k(\mathcal{O}_Y, f_*\mathcal{O}_X) \ar[lld] \\
& T\Deformationcategory_{X \to Y} \ar[r] &
T(\Deformationcategory_X \times \Deformationcategory_Y) \ar[r] &
\text{Ext}^1_{\mathcal{O}_X}(Lf^*\NL_{Y/k}, \mathcal{O}_X)
}
$$
\end{lemma}

\begin{proof}
変形圏の明らかな写像
$\Deformationcategory_{X \to Y} \to
\Deformationcategory_X \times \Deformationcategory_Y$
は、補題の完全列に現れる射のうち2つを与える。
$\text{Inf}(\Deformationcategory_{X \to Y})$ は、$X \to Y$ の
$k[\epsilon]$ への自明変形
$$
f' :  X' = X \times_{\Spec(k)} \Spec(k[\epsilon])
\xrightarrow{f \times \text{id}}
Y' = Y \times_{\Spec(k)} \Spec(k[\epsilon])
$$
の自己同型で、$\epsilon$ を法として恒等写像となるもの全体の集合である。
これは明らかに、$f'$ と両立する $X$ と $Y$ の無限小自己同型の対
$(\alpha, \beta) \in
\text{Inf}(\Deformationcategory_X \times \Deformationcategory_Y)$
、すなわち $f' \circ \alpha = \beta \circ f'$ を満たす対にほかならない。
『変形理論』の補題 \ref{defos-lemma-huge-diagram-ringed-spaces} により、
任意の対 $(\alpha, \beta)$ に対して、射 $f' : X' \to Y'$ と射
$\beta^{-1} \circ f' \circ \alpha : X' \to Y'$ との差は
次の空間の元を定める。
$$
\text{Der}_k(\mathcal{O}_Y, f_*\mathcal{O}_X) =
\Hom_{\mathcal{O}_Y}(\Omega_{Y/k}, f_*\mathcal{O}_X)
$$
この等式は『射の詳論』の補題
\ref{more-morphisms-lemma-netherlander-quasi-coherent} による。
これによって上段の最後の水平射が定まり、最初の2か所での完全性が示される。
写像
$$
\text{Der}_k(\mathcal{O}_Y, f_*\mathcal{O}_X)
\to
T\Deformationcategory_{X \to Y}
$$
については、『変形理論』の補題
\ref{defos-lemma-huge-diagram-ringed-spaces} を用い、始域の元を
$\epsilon$ を法として $f$ に等しい $\Spec(k[\epsilon])$ 上の射
$f_\epsilon : X' \to Y'$ と解釈する。$f_\epsilon$ を
$T\Deformationcategory_{X \to Y}$ における
$(f_\epsilon : X' \to Y')$ の同型類に送る。
$(f_\epsilon : X' \to Y')$ が自明変形 $(f' : X' \to Y')$ と同型であるのは、
ある対 $(\alpha, \beta)$ に対して
$f_\epsilon  = \beta^{-1} \circ f \circ \alpha$ となるとき、かつそのときに限る。
% P11-E0053
これにより第3項での完全性が従う。明らかに、ある一次変形
$(f_\epsilon : X_\epsilon \to Y_\epsilon)$
が $T(\Deformationcategory_X \times \Deformationcategory_Y)$ で零に写るならば、
同型 $X' \to X_\epsilon$ および $Y' \to Y_\epsilon$ を選べるので、
この元は南西向きの射の像に属する。したがって第4項でも完全である。
最後に、$X$ と $Y$ の2つの一次変形 $X_\epsilon$、$Y_\epsilon$ が与えられると、
次の空間に障害が存在する。
$$
ob(X_\epsilon, Y_\epsilon) \in
\text{Ext}^1_{\mathcal{O}_X}(Lf^*\NL_{Y/k}, \mathcal{O}_X)
$$
これは $f : X \to Y$ が $X_\epsilon \to Y_\epsilon$ に持ち上がるとき、
かつそのときに限り消える。『変形理論』の補題
\ref{defos-lemma-huge-diagram-ringed-spaces} を参照せよ。これで証明が完了する。
\end{proof}

\begin{lemma}
\label{examples-defos-lemma-proper-schemes-morphisms-TI}
補題 \ref{examples-defos-lemma-schemes-morphisms-TI} において $X$ と $Y$ がともに
$k$ 上固有ならば、$\text{Inf}(\Deformationcategory_{X \to Y})$ と
$T\Deformationcategory_{X \to Y}$ は有限次元である。
\end{lemma}

\begin{proof}
省略する。ヒント：補題 \ref{examples-defos-lemma-proper-schemes-TI} と同様に論じ、
本補題の完全列を用いよ。
\end{proof}

\noindent
例 \ref{examples-defos-example-schemes-morphisms} において $X \to Y$ が
$k$ 上の固有スキームの射ならば、$\Deformationcategory_{X \to Y}$ は
$\mathcal{C}_\Lambda$ 上の関手のなす滑らかなプロ表現可能グルーポイドによる
表示をもち、したがって特に（極小）半普遍的形式対象をもつ。これは補題
\ref{examples-defos-lemma-schemes-morphisms-RS}、
\ref{examples-defos-lemma-proper-schemes-morphisms-TI}
および節 \ref{examples-defos-section-general} の一般的な議論から従う。

\begin{lemma}
\label{examples-defos-lemma-schemes-morphisms-hull}
例 \ref{examples-defos-example-schemes-morphisms} において $X \to Y$ を
固有 $k$ スキームの射とする。$\Lambda$ を剰余体 $k$ をもつ完備局所環と仮定する
（古典的な場合）。このとき関手
$$
F : \mathcal{C}_\Lambda \longrightarrow \textit{Sets},\quad
A \longmapsto \Ob(\Deformationcategory_{X \to Y}(A))/\cong
$$
は、対象の同型類を対応させる関手であり、射影極限的包をもつ。さらに
$\text{Der}_k(\mathcal{O}_X, \mathcal{O}_X) =
\text{Der}_k(\mathcal{O}_Y, \mathcal{O}_Y) = 0$ ならば、
$F$ はプロ表現可能である。
\end{lemma}

\begin{proof}
射影極限的包の存在は、補題 \ref{examples-defos-lemma-schemes-morphisms-RS}、
\ref{examples-defos-lemma-proper-schemes-morphisms-TI}、『形式的変形理論』の補題
\ref{formal-defos-lemma-RS-implies-S1-S2}、および注意
\ref{formal-defos-remark-compose-minimal-into-iso-classes} から直ちに従う。

\medskip\noindent
$\text{Der}_k(\mathcal{O}_X, \mathcal{O}_X) =
\text{Der}_k(\mathcal{O}_Y, \mathcal{O}_Y) = 0$ と仮定する。このとき
補題 \ref{examples-defos-lemma-schemes-morphisms-TI} の完全列と補題
\ref{examples-defos-lemma-schemes-TI} を合わせると、
$\text{Inf}(\Deformationcategory_{X \to Y}) = 0$ が分かる。
したがって『形式的変形理論』の補題
\ref{formal-defos-lemma-infdef-trivial} により、
$\Deformationcategory_{X \to Y}$ と $F$ は同値である。ゆえに $F$ は
有限次元接空間をもつ変形関手であり
（$\Deformationcategory_{X \to Y}$ は変形圏である）、
『形式的変形理論』の定理
\ref{formal-defos-theorem-Schlessinger-prorepresentability} を適用できる。
\end{proof}

\begin{lemma}
\label{examples-defos-lemma-schemes-morphisms-smooth-to-base}
\begin{reference}
これについては \cite[Section 5.3]{Ravi-Murphys-Law} および
\cite[Theorem 3.3]{Ran-deformations} で論じられている。
\end{reference}
例 \ref{examples-defos-example-schemes} において、$f : X \to Y$ を $k$ 上のスキームの射とする。
$f_*\mathcal{O}_X = \mathcal{O}_Y$ かつ $R^1f_*\mathcal{O}_X = 0$ ならば、
変形圏の射
$$
\Deformationcategory_{X \to Y} \to \Deformationcategory_X
$$
は同値である。
\end{lemma}

\begin{proof}
本補題の忘却関手の擬逆を構成する。$(A, U)$ を
$\Deformationcategory_X$ の対象とする。与えられた写像 $X \to U$ は
有限位数の厚化であるから、これを用いて $U$ と $X$ の基礎位相空間を
同一視できる。『射の詳論』の節
\ref{more-morphisms-section-thickenings} を参照せよ。したがって
$\mathcal{O}_U$ を $X$ 上の $A$ 代数の層とみなすことができ、以下そうする。
さらに $U \to \Spec(A)$ が平坦であることは、$\mathcal{O}_U$ が
$A$ 加群の層として平坦であることを意味する。特に、フィルトレーション
$$
0 = \mathfrak m_A^n\mathcal{O}_U \subset
\mathfrak m_A^{n - 1}\mathcal{O}_U \subset \ldots \subset
\mathfrak m_A^2\mathcal{O}_U \subset
\mathfrak m_A\mathcal{O}_U \subset \mathcal{O}_U
$$
をもち、その各商は平坦性により
$\mathcal{O}_X \otimes_k \mathfrak m_A^i/\mathfrak m_A^{i + 1}$
に等しい。『射の詳論』の補題 \ref{more-morphisms-lemma-deform}、
または、より一般的な『変形理論』の補題
\ref{defos-lemma-deform-module} を参照せよ。次のようにおく。
$$
\mathcal{O}_V = f_*\mathcal{O}_U
$$
これは $Y$ 上の $A$ 代数の層とみなす。$R^1f_*\mathcal{O}_X = 0$ なので、
上記の記述から、すべての $i$ に対して
$R^1f_*(\mathfrak m_A^i\mathcal{O}_U/\mathfrak m_A^{i + 1}\mathcal{O}_U) = 0$
であることが分かる。したがって、すべての $i$ に対して列
$$
0 \to
(f_*\mathcal{O}_X) \otimes_k \mathfrak m_A^i/\mathfrak m_A^{i + 1} \to
f_*(\mathcal{O}_U/\mathfrak m_A^{i + 1}\mathcal{O}_U) \to
f_*(\mathcal{O}_U/\mathfrak m_A^i\mathcal{O}_U) \to 0
$$
は完全である。上記の参照結果を逆向きに用い（さらに帰納法を用いると）、
$\mathcal{O}_V$ は
$\mathcal{O}_V/\mathfrak m_A\mathcal{O}_V = \mathcal{O}_Y$ を満たす
$A$ 代数の平坦な層であることが分かる。『射の詳論』の補題
\ref{more-morphisms-lemma-first-order-thickening}
により、$(Y, \mathcal{O}_V)$ はスキームである。これを $V$ と呼ぶ。
等式 $\mathcal{O}_V = f_*\mathcal{O}_U$ は環付き空間の射
$U \to V$ を定め、これがスキームの射であることは容易に分かる。
すでに示した平坦性と合わせて、証明が完了する。
\end{proof}







\section{代数空間}
\label{examples-defos-section-algebraic-spaces}

\noindent
代数空間の変形理論を扱う。

\begin{example}[代数空間]
\label{examples-defos-example-spaces}
圏 $\mathcal{F}$ を次のように定義する。
\begin{enumerate}
\item 対象は、$\mathcal{C}_\Lambda$ の対象 $A$ と $A$ 上平坦な代数空間 $X$
からなる対 $(A, X)$ である。
\item 射 $(f, g) : (B, Y) \to (A, X)$ は、$\mathcal{C}_\Lambda$ における
射 $f : B \to A$ と、次が代数空間のデカルト可換図式となるような
$\Lambda$ 上の代数空間の射 $g : X \to Y$ からなる。
$$
\xymatrix{
X \ar[r]_g \ar[d] & Y \ar[d] \\
\Spec(A) \ar[r]^f & \Spec(B)
}
$$
\end{enumerate}
関手 $p : \mathcal{F} \to \mathcal{C}_\Lambda$ は $(A, X)$ を $A$ に、
$(f, g)$ を $f$ に送る。$p$ がグルーポイドをファイバーとする
余ファイバー圏であることは明らかである。$k$ 上の代数空間 $X$ に対し、
$x_0 = (k, X)$ を対応する $\mathcal{F}(k)$ の対象とする。次のようにおく。
$$
\Deformationcategory_X = \mathcal{F}_{x_0}
$$
\end{example}

\begin{lemma}
\label{examples-defos-lemma-spaces-RS}
例 \ref{examples-defos-example-spaces} は Rim--Schlessinger 条件 (RS) を満たす。
特に、任意の $k$ 上の代数空間 $X$ に対して
$\Deformationcategory_X$ は変形圏である。
\end{lemma}

\begin{proof}
$A_1 \to A$ と $A_2 \to A$ を $\mathcal{C}_\Lambda$ の射とし、
$A_2 \to A$ は全射であると仮定する。『形式的変形理論』の補題
\ref{formal-defos-lemma-RS-2-categorical}
により、関手
$\mathcal{F}(A_1 \times_A A_2) \to
\mathcal{F}(A_1) \times_{\mathcal{F}(A)} \mathcal{F}(A_2)$
が圏同値であることを示せば十分である。図式
$$
\xymatrix{
\Spec(A) \ar[r] \ar[d] & \Spec(A_2) \ar[d] \\
\Spec(A_1) \ar[r] &
\Spec(A_1 \times_A A_2)
}
$$
は『空間の押し出し』の補題
\ref{spaces-pushouts-lemma-pushout-along-thickening} にいう押し出し図式である。
したがって本補題は『空間の押し出し』の補題
\ref{spaces-pushouts-lemma-equivalence-categories-spaces-pushout-flat}
の特別な場合である。
\end{proof}

\begin{lemma}
\label{examples-defos-lemma-spaces-TI}
例 \ref{examples-defos-example-spaces} において、$X$ を $k$ 上の代数空間とする。このとき
$$
\text{Inf}(\Deformationcategory_X) =
\text{Ext}^0_{\mathcal{O}_X}(\NL_{X/k}, \mathcal{O}_X) =
\Hom_{\mathcal{O}_X}(\Omega_{X/k}, \mathcal{O}_X) =
\text{Der}_k(\mathcal{O}_X, \mathcal{O}_X)
$$
かつ
$$
T\Deformationcategory_X =
\text{Ext}^1_{\mathcal{O}_X}(\NL_{X/k}, \mathcal{O}_X)
$$
\end{lemma}

\begin{proof}
$\text{Inf}(\Deformationcategory_X)$ は、$X$ の $k[\epsilon]$ への自明変形
$X' = X \times_{\Spec(k)} \Spec(k[\epsilon])$ の自己同型で、
$\epsilon$ を法として恒等写像となるもの全体の集合であることを思い出そう。
『変形理論』の補題 \ref{defos-lemma-deform-spaces} により、これは
$\text{Ext}^0_{\mathcal{O}_X}(\NL_{X/k}, \mathcal{O}_X)$ に等しい。
等式 $\text{Ext}^0_{\mathcal{O}_X}(\NL_{X/k}, \mathcal{O}_X) =
\Hom_{\mathcal{O}_X}(\Omega_{X/k}, \mathcal{O}_X)$ は、
『空間の射の詳論』の補題
\ref{spaces-more-morphisms-lemma-netherlander-quasi-coherent} から従う。
また、等式
$\Hom_{\mathcal{O}_X}(\Omega_{X/k}, \mathcal{O}_X) =
\text{Der}_k(\mathcal{O}_X, \mathcal{O}_X)$
は『空間の射の詳論』の定義
\ref{spaces-more-morphisms-definition-sheaf-differentials} と
『サイト上の加群』の定義
\ref{sites-modules-definition-module-differentials} から従う。

\medskip\noindent
$T_{x_0}\Deformationcategory_X$ は、$X$ の $k[\epsilon]$ への平坦変形
$X'$ の同型類全体の集合、より正確には
$\Deformationcategory_X(k[\epsilon])$ の同型類全体の集合であることを思い出そう。
したがって補題の第2の主張は、『変形理論』の補題
\ref{defos-lemma-deform-spaces} から従う。
\end{proof}

\begin{lemma}
\label{examples-defos-lemma-proper-spaces-TI}
補題 \ref{examples-defos-lemma-spaces-TI} において $X$ が $k$ 上固有ならば、
$\text{Inf}(\Deformationcategory_X)$ と $T\Deformationcategory_X$ は
有限次元である。
\end{lemma}

\begin{proof}
同補題により、
$\Ext^1_{\mathcal{O}_X}(\NL_{X/k}, \mathcal{O}_X)$ と
$\Ext^0_{\mathcal{O}_X}(\NL_{X/k}, \mathcal{O}_X)$ が有限次元であることを
示せばよい。『空間の射の詳論』の補題
\ref{spaces-more-morphisms-lemma-netherlander-fp}
と $X$ がネーターであることから、$\NL_{X/k}$ のコホモロジー層は
連接であり、次数 $0$ と $-1$ を除いて零であることが分かる。
『空間の導来圏』の補題 \ref{spaces-perfect-lemma-ext-finite} により、
上に表示した $\Ext$ 群は有限次元 $k$ ベクトル空間である。これで証明が完了する。
\end{proof}

\noindent
例 \ref{examples-defos-example-spaces} において $X$ が $k$ 上の固有代数空間ならば、
$\Deformationcategory_X$ は $\mathcal{C}_\Lambda$ 上の関手のなす
滑らかなプロ表現可能グルーポイドによる表示をもち、したがって特に
（極小）半普遍的形式対象をもつ。これは補題 \ref{examples-defos-lemma-spaces-RS}、
\ref{examples-defos-lemma-proper-spaces-TI}
および節 \ref{examples-defos-section-general} の一般的な議論から従う。

\begin{lemma}
\label{examples-defos-lemma-spaces-hull}
例 \ref{examples-defos-example-spaces} において $X$ を $k$ 上の固有代数空間とする。
$\Lambda$ を剰余体 $k$ をもつ完備局所環と仮定する（古典的な場合）。
このとき関手
$$
F : \mathcal{C}_\Lambda \longrightarrow \textit{Sets},\quad
A \longmapsto \Ob(\Deformationcategory_X(A))/\cong
$$
は、対象の同型類を対応させる関手であり、射影極限的包をもつ。さらに
$\text{Der}_k(\mathcal{O}_X, \mathcal{O}_X) = 0$ ならば、
$F$ はプロ表現可能である。
\end{lemma}

\begin{proof}
射影極限的包の存在は、補題 \ref{examples-defos-lemma-spaces-RS}、
\ref{examples-defos-lemma-proper-spaces-TI}、『形式的変形理論』の補題
\ref{formal-defos-lemma-RS-implies-S1-S2}、および注意
\ref{formal-defos-remark-compose-minimal-into-iso-classes} から直ちに従う。

\medskip\noindent
$\text{Der}_k(\mathcal{O}_X, \mathcal{O}_X) = 0$ と仮定する。
このとき『形式的変形理論』の補題
\ref{formal-defos-lemma-infdef-trivial} により、
$\Deformationcategory_X$ と $F$ は同値である。したがって $F$ は
有限次元接空間をもつ変形関手であり（$\Deformationcategory_X$ は変形圏である）、
『形式的変形理論』の定理
\ref{formal-defos-theorem-Schlessinger-prorepresentability} を適用できる。
\end{proof}





\section{完備化の変形}
\label{examples-defos-section-compare}

\noindent
本節では、代数とその完備化が定める変形問題を比較する。
まず「持ち上げ可能性」について論じる。

\begin{lemma}
\label{examples-defos-lemma-lift-equivalence-module-derived}
$A' \to A$ を冪零核をもつ環の全射とし、$A' \to P'$ を平坦な環準同型とする。
$P = P' \otimes_{A'} A$ とおく。$M$ を $A$ 上平坦な $P$ 加群とする。
このとき次は同値である。
\begin{enumerate}
\item $M' \otimes_{P'} P = M$ を満たす $A'$ 上平坦な $P'$ 加群 $M'$ が存在する。
\item $K' \otimes_{P'}^\mathbf{L} P = M$ を満たす対象 $K' \in D^-(P')$ が存在する。
\end{enumerate}
\end{lemma}

\begin{proof}
(1) のような $M'$ が存在すると仮定する。このとき
$$
M = M' \otimes_P P' = M' \otimes_{A'} A =
M' \otimes_A^\mathbf{L} A' = M' \otimes_{P'}^\mathbf{L} P
$$
% P11-E0054
最初の2つの等式は明らかであり、第3の等式は $M'$ が $A'$ 上平坦であることから、
第4の等式は『代数の詳論』の補題
\ref{more-algebra-lemma-base-change-comparison} から従う。したがって (2) が成り立つ。
逆に、(2) のような $K'$ が存在すると仮定する。$M$ は非零であると仮定してよく、
以下そうする。$H^t(K')$ が非零となる最大の整数を $t$ とする
（$M$ が非零なので存在する）。$t > 0$ ならば
$H^t(K') \otimes_{P'} P = H^t(K' \otimes_{P'}^\mathbf{L} P)$ は零である。
$P' \to P$ の核は冪零なので、これは中山の補題により
$H^t(K') = 0$ を意味し、矛盾である。したがって $t = 0$ である
（$t < 0$ の場合も不可能である）。このとき $M' = H^0(K')$ は
$M = M' \otimes_{P'} P$ を満たす $P'$ 加群であり、Tor のスペクトル系列から
単射
$$
\text{Tor}_1^{P'}(M', P) \to H^{-1}(M' \otimes_{P'}^\mathbf{L} P) = 0
$$
を得る。上で挙げた導来基底変換に関する結果により、
$0 = \text{Tor}_1^{P'}(M', P) = \text{Tor}_1^{A'}(M', A)$ である。
『代数』の補題 \ref{algebra-lemma-what-does-it-mean} により、
$M'$ は $A'$ 上平坦である。
\end{proof}

\begin{lemma}
\label{examples-defos-lemma-lift-equivalence-module}
ネーター環の可換図式
$$
\xymatrix{
A' \ar[d] \ar[r] &
P' \ar[d] \ar[r] &
Q' \ar[d] \\
A \ar[r] &
P \ar[r] &
Q
}
$$
を考える。両正方形はデカルトであり、水平射は平坦、垂直射は全射で
その核は冪零であるとする。$J' \subset P'$ を
$P'/J' = Q'/J'Q'$ を満たすイデアルとする。$M$ を $A$ 上平坦な
$P$ 加群とする。すべての $g \in J'$ に対し、$M_g$ を持ち上げる
$A'$ 上平坦な $(P')_g$ 加群が存在すると仮定する。このとき次は同値である。
\begin{enumerate}
\item $M$ は $P'$ 加群への $A'$ 上平坦な持ち上げをもつ。
\item $M \otimes_P Q$ は $Q'$ 加群への $A'$ 上平坦な持ち上げをもつ。
\end{enumerate}
\end{lemma}

\begin{proof}
$I = \Ker(A' \to A)$ とおく。$I^n = 0$ を満たす整数 $n > 1$ に関する
帰納法により、$I$ が平方零イデアルの場合に帰着する。詳細は省略する。
$M$ の持ち上げ可能性という条件を、補題
\ref{examples-defos-lemma-lift-equivalence-module-derived} のような $D^-(P')$ の対象を
見つける問題に読み替える。その障害は元
$$
\omega(M) \in \text{Ext}^2_P(M, M \otimes_P^\mathbf{L} IP) =
\text{Ext}^2_P(M, M \otimes_P IP)
$$
であり、『変形理論』の補題
\ref{defos-lemma-canonical-class-algebra} で構成されている。
上の表示式の等号が成り立つのは、$M$ と $P$ が $A$ 上平坦なので
$M \otimes_P^\mathbf{L} IP = M \otimes_P IP$
だからである\footnote{自由 $A$ 加群による分解
$F_\bullet \to I$ を選ぶ。$A \to P$ は平坦なので、
$P \otimes_A F_\bullet$ は $IP$ の自由分解である。したがって
$M \otimes_P^\mathbf{L} IP$ は
$M \otimes_P P \otimes_A F_\bullet = M \otimes_A F_\bullet$ で表される。
$M$ は $A$ 上平坦なので、これは次数 $0$ にしかコホモロジーをもたない。}。
同様に、$M \otimes_P Q$ を持ち上げる障害は元
$$
\omega(M \otimes_P Q) \in
\text{Ext}^2_Q(M \otimes_P Q, (M \otimes_P Q) \otimes_Q IQ)
$$
である。これは『変形理論』の補題
\ref{defos-lemma-canonical-class-algebra} による構成 $\omega(-)$ の関手性のもとでの
$\omega(M)$ の像である。『代数の詳論』の補題
\ref{more-algebra-lemma-base-change-RHom} により
$$
\text{Ext}^2_Q(M \otimes_P Q, (M \otimes_P Q) \otimes_Q IQ) =
\text{Ext}^2_P(M, M \otimes_P IP) \otimes_P Q
$$
を得る。ここでは $P$ がネーターであり $M$ が有限であることを用いた。
% P11-E0055
$P' \to Q'$ に関する仮定から、$P$ 加群 $E$ に対する写像
$E \to E \otimes_P Q$ は $J'$ 冪捩れ部分上で全単射である。
% P11-E0057
『代数の詳論』の補題
\ref{more-algebra-lemma-neighbourhood-equivalence}.
を参照せよ。したがって $\omega(M)$ が $J'$ 冪捩れであることを示せば十分である。
言い換えれば、すべての $g \in J'$ に対し $\omega(M)$ が
$$
\text{Ext}^2_P(M, M \otimes_P IP)_g =
\text{Ext}^2_{P_g}(M_g, M_g \otimes_{P_g} IP_g)
$$
で消えることを示せば十分である。しかし、$\omega(M)$ の構成は基底変換と
両立するので、$M_g$ が持ち上げをもつという仮定から、これは成り立つ
（もちろん、同値の連鎖全体を再び用いる必要がある）。
% P11-E0056
\end{proof}

\begin{lemma}
\label{examples-defos-lemma-lift-equivalence}
$A' \to A$ を冪零核をもつネーター環の全射とし、$A \to B$ を
有限型の平坦な環準同型とする。$\mathfrak b \subset B$ を、
$\Spec(B) \to \Spec(A)$ が $V(\mathfrak b)$ の補集合上シントミックとなる
イデアルとする。このとき $B$ が $A'$ への平坦な持ち上げをもつことと、
$\mathfrak b$ 進完備化 $B^\wedge$ が $A'$ への平坦な持ち上げをもつことは同値である。
\end{lemma}

\begin{proof}
$A$ 代数の全射 $P = A[x_1, \ldots, x_n] \to B$ を選ぶ。
$\mathfrak p \subset P$ を $\mathfrak b$ の逆像とする。
$P' = A'[x_1, \ldots, x_n]$ とおき、$\mathfrak p' \subset P'$ で
$\mathfrak p$ の逆像を表す（もちろん、ここで $\mathfrak p$ と
$\mathfrak p'$ は素イデアルを表すものではない）。それぞれの完備化を
$P^\wedge$ および $(P')^\wedge$ と書く。

\medskip\noindent
$A' \to B'$ を $A \to B$ の平坦な持ち上げとする。すなわち、
$A' \to B'$ は平坦で、$A$ 代数の同型
$B = B' \otimes_{A'} A$ が存在するとする。このとき、与えられた全射
$P \to B$ を持ち上げる $A'$ 代数準同型 $P' \to B'$ を選べる。
中山の補題（『代数』の補題 \ref{algebra-lemma-NAK}）により、
$B'$ は $P'$ の商である。特に、$A$ 上平坦な $P$ 加群としての $B$ を
持ち上げる、$A'$ 上平坦な $P'$ 加群構造を $B'$ に入れられる。
逆に、$B$ を $A'$ 上平坦な $P'$ 加群 $M'$ に持ち上げられるならば、
$M'$ は巡回加群 $M' \cong P'/J'$ である（ここでも中山の補題を用いる）。
$B' = P'/J'$ とおけば、代数としての $B$ の平坦な持ち上げが得られる。

\medskip\noindent
$C = B^\wedge$、$\mathfrak c = \mathfrak bC$ とおく。
$A' \to C'$ を $A \to C$ の平坦な持ち上げとする。このとき $C'$ は、
$\mathfrak c$ の逆像 $\mathfrak c'$ に関して完備である
（『代数』の補題 \ref{algebra-lemma-complete-modulo-nilpotent}）。
$A$ 代数準同型 $P \to C$ を持ち上げる $A'$ 代数準同型 $P' \to C'$ を選ぶ。
これらの準同型は完備化を経由し、全射 $P^\wedge \to C$ および
$(P')^\wedge \to C'$ を与える（後者には再び中山の補題を用いる）。
特に、$A$ 上平坦な $P^\wedge$ 加群としての $C$ を持ち上げる、
$A'$ 上平坦な $(P')^\wedge$ 加群構造を $C'$ に入れられる。
逆に、$C$ を $A'$ 上平坦な $(P')^\wedge$ 加群 $N'$ に持ち上げられるならば、
$N'$ は巡回加群 $N' \cong (P')^\wedge/\tilde J$ である
（ここでも中山の補題を用いる）。$C' = (P')^\wedge/\tilde J$ とおけば、
代数としての $C$ の平坦な持ち上げが得られる。

\medskip\noindent
$P' \to (P')^\wedge$ は平坦な環準同型であり、同型
$P'/\mathfrak p' = (P')^\wedge/\mathfrak p'(P')^\wedge$ を誘導する。
したがって、すべての $g \in \mathfrak p'$ に対して $B_g$ が
$A'$ 上平坦な $P'_g$ 加群に持ち上がることを示せば、本補題は補題
\ref{examples-defos-lemma-lift-equivalence-module} から従う。しかし環準同型
$A \to B_g$ はシントミックなので、『環準同型の平滑化』の命題
\ref{smoothing-proposition-lift-smooth} により $A'$ 上平坦な代数 $B'$ に持ち上がる。
$A' \to P'_g$ は滑らかなので、前と同様に $P_g \to B_g$ を全射
$P'_g \to B'$ に持ち上げられ、求めるものを得る。
\end{proof}

\noindent
記法。$A \to B$ を環準同型とし、$N$ を $B$ 加群とする。
$N$ が $C$ の平方零イデアルとなる $A$ 代数の拡大
$$
0 \to N \to C \to B \to 0
$$
の同型類全体の集合を $\text{Exal}_A(B, N)$ と書く。もう1つのこのような拡大
$0 \to N \to C' \to B \to 0$ が与えられたとき、同型とは、次を可換にする
$A$ 代数の同型 $C \to C'$ のことである。
% P11-E0059
$$
\xymatrix{
0 \ar[r] &
N \ar[r] \ar[d]_{\text{id}} &
C \ar[r] \ar[d] &
B \ar[r] \ar[d]_{\text{id}} & 0 \\
0 \ar[r] &
N \ar[r] &
C' \ar[r] &
B \ar[r] & 0
}
$$
対応 $N \mapsto \text{Exal}_A(B, N)$ は積を積に移す関手である。
したがってこれは加法関手であり、$\text{Exal}_A(B, N)$ は自然な
$B$ 加群構造をもつ。実際、『変形理論』の補題
\ref{defos-lemma-choices} により
$\text{Exal}_A(B, N) = \text{Ext}^1_B(\NL_{B/A}, N)$ である。

\begin{lemma}
\label{examples-defos-lemma-first-order-completion}
$k$ を体、$B$ を有限型 $k$ 代数とする。$J \subset B$ を、
$\Spec(B) \to \Spec(k)$ が $V(J)$ の補集合上滑らかとなるイデアルとする。
$N$ を有限 $B$ 加群とする。このとき標準的な全単射
$$
\text{Exal}_k(B, N) \to \text{Exal}_k(B^\wedge, N^\wedge)
$$
が存在する。ここで $B^\wedge$ と $N^\wedge$ は $J$ 進完備化である。
\end{lemma}

\begin{proof}
この写像は完備化によって与えられる。すなわち、$\text{Exal}_k(B, N)$ の元
$0 \to N \to C \to B \to 0$ を、$J$ の逆像に関する $C$ の完備化
$C^\wedge$ に送る。補題 \ref{examples-defos-lemma-completion} の証明と比較せよ。

\medskip\noindent
$k \to B$ は有限表示なので、複体 $\NL_{B/k}$ は、$N^i$ が有限 $B$ 加群である
複体 $N^{-1} \to N^0$ によって表せる。『代数』の節
\ref{algebra-section-netherlander}、特に補題
\ref{algebra-lemma-NL-homotopy} を参照せよ。$B$ はネーターなので、
これは $\NL_{B/k}$ が擬連接であることを意味する。$g \in J$ に対して
$k$ 代数 $B_g$ は滑らかであり、したがって
$(\NL_{B/k})_g = \NL_{B_g/k}$ は次数 $0$ に置かれた有限射影 $B$ 加群と
擬同型である。
% P11-E0058
ゆえに任意の $B$ 加群 $N$ と $i \geq 1$ に対して
$\text{Ext}^i_B(\NL_{B/k}, N)_g = 0$ である。『代数の詳論』の補題
\ref{more-algebra-lemma-ext-annihilated-into} により、
$$
\text{Ext}^1_B(\NL_{B/k}, N) \longrightarrow
\lim_n \text{Ext}^1_B(\NL_{B/k}, N/J^n N)
$$
は任意の有限 $B$ 加群 $N$ に対して同型である。

\medskip\noindent
写像の単射性を示す。$0 \to N \to C \to B \to 0$ を
$\text{Exal}_k(B, N)$ の元で、$\text{Exal}_k(B^\wedge, N^\wedge)$ で
零に写るものとする。分裂 $C^\wedge = B^\wedge \oplus N^\wedge$ を選ぶ。
このとき誘導される写像 $C \to C^\wedge \to N^\wedge$ は、すべての $n$ に対して
写像 $C \to N/J^nN$ を与える。したがって、この元はすべての $n$ に対して写像
$$
\text{Ext}^1_B(\NL_{B/k}, N) \to
\text{Ext}^1_B(\NL_{B/k}, N/J^n N)
$$
の核に属する。前段落により、この元は零である。

\medskip\noindent
写像の全射性を示す。$0 \to N^\wedge \to C' \to B^\wedge \to 0$ を
$\text{Exal}_k(B^\wedge, N^\wedge)$ の元とする。$B \to B^\wedge$ によって
引き戻すと、$\text{Exal}_k(B, N^\wedge)$ の元
$0 \to N^\wedge \to C'' \to B \to 0$ を得る。次の等式が成り立つ。
$$
\text{Ext}^1_B(\NL_{B/k}, N^\wedge) =
\text{Ext}^1_B(\NL_{B/k}, N) \otimes_B B^\wedge =
\text{Ext}^1_B(\NL_{B/k}, N)
$$
第1の等式は、$N^\wedge = N \otimes_B B^\wedge$
（『代数』の補題 \ref{algebra-lemma-completion-tensor}）および
『代数の詳論』の補題
\ref{more-algebra-lemma-pseudo-coherence-and-ext} による。
第2の等式は、$\text{Ext}^1_B(\NL_{B/k}, N)$ が $J$ 冪捩れであること
（上記参照）、$B \to B^\wedge$ が平坦で同型
$B/J \to B^\wedge/JB^\wedge$ を誘導すること、および『代数の詳論』の補題
\ref{more-algebra-lemma-neighbourhood-equivalence} による。したがって、
$\text{Exal}_k(B, N^\wedge)$ で $C''$ に写る
$C \in \text{Exal}_k(B, N)$ を見つけられる。ゆえに
$$
0 \to N^\wedge \to C' \to B^\wedge \to 0
\quad\text{and}\quad
0 \to N^\wedge \to C^\wedge \to B^\wedge \to 0
$$
は $\text{Exal}_k(B^\wedge, N^\wedge)$ の2つの元であり、
$\text{Exal}_k(B, N^\wedge)$ の同じ元に写る。差を取ると、
$\text{Exal}_k(B^\wedge, N^\wedge)$ の元
$0 \to N^\wedge \to C' \to B^\wedge \to 0$ で、
$\text{Exal}_k(B, N^\wedge)$ における像が零となるものを得る。
これは次の図式が存在することを意味する。
$$
\xymatrix{
0 \ar[r] &
N^\wedge \ar[r] &
C' \ar[r] &
B^\wedge \ar[r] & 0 \\
& & B \ar[u]^\sigma \ar[ru]
}
$$
$J' \subset C'$ を $JB^\wedge \subset B^\wedge$ の逆像とする。
証明を終えるには、$\sigma$ が $B$ の $J$ 進位相と $C'$ の $J'$ 進位相に関して
連続であり、『代数』の補題
\ref{algebra-lemma-complete-modulo-nilpotent} により $C'$ が $J'$ 進完備であることに
注意すればよい（ここでは $C'$ がネーターであることも用いる。細部は省略する）。
実際、これは $\sigma$ が完備化 $B^\wedge$ を経由し、
$\text{Exal}_k(B^\wedge, N^\wedge)$ において $C' = 0$ となることを意味する。
\end{proof}

\begin{lemma}
\label{examples-defos-lemma-smooth-completion}
例 \ref{examples-defos-example-rings} において $P$ を $k$ 代数とし、$J \subset P$ を
イデアルとする。$J$ 進完備化を $P^\wedge$ と書く。もし
\begin{enumerate}
\item $k \to P$ が有限型であり、
\item $\Spec(P) \to \Spec(k)$ が $V(J)$ の補集合上滑らかである
\end{enumerate}
ならば、補題 \ref{examples-defos-lemma-completion} の変形圏の間の関手
$$
\Deformationcategory_P \longrightarrow \Deformationcategory_{P^\wedge}
$$
は滑らかであり、接空間上の同型を誘導する。
\end{lemma}

\begin{proof}
補題 \ref{examples-defos-lemma-rings-RS} により、$\Deformationcategory_P$ と
$\Deformationcategory_{P^\wedge}$ は変形圏である。したがって、
本関手が接空間を同一視し、持ち上げ可能性が対応することを確認すれば十分である。
『形式的変形理論』の補題
\ref{formal-defos-lemma-easy-check-smooth} を参照せよ。持ち上げ可能性に関する性質は
補題 \ref{examples-defos-lemma-lift-equivalence} で示され、接空間上の同型は補題
\ref{examples-defos-lemma-first-order-completion} における $N = B$ の特別な場合である。
\end{proof}



\section{局所化の変形}
\label{examples-defos-section-compare-localization}

\noindent
本節では、代数とその乗法的部分集合における局所化が定める変形問題を比較する。
まず「持ち上げ可能性」について論じる。

\begin{lemma}
\label{examples-defos-lemma-lift-equivalence-localization}
$A' \to A$ を冪零核をもつネーター環の全射とし、$A \to B$ を
有限型の平坦な環準同型とする。$S \subset B$ を、
$\Spec(B) \to \Spec(A)$ が $\mathfrak q$ においてシントミックでなければ
$S \cap \mathfrak q = \emptyset$ となる乗法的部分集合とする。このとき
$B$ が $A'$ への平坦な持ち上げをもつことと、$S^{-1}B$ が
$A'$ への平坦な持ち上げをもつことは同値である。
\end{lemma}

\begin{proof}
この証明は補題 \ref{examples-defos-lemma-lift-equivalence} の証明と同じだが、より易しい。
読者には証明を飛ばすことを勧める。
$A$ 代数の全射 $P = A[x_1, \ldots, x_n] \to B$ を選ぶ。
$S_P \subset P$ を $S$ の逆像とする。$P' = A'[x_1, \ldots, x_n]$ とおき、
$S_{P'} \subset P'$ で $S_P$ の逆像を表す。

\medskip\noindent
$A' \to B'$ を $A \to B$ の平坦な持ち上げとする。すなわち、
$A' \to B'$ は平坦で、$A$ 代数の同型
$B = B' \otimes_{A'} A$ が存在するとする。このとき、与えられた全射
$P \to B$ を持ち上げる $A'$ 代数準同型 $P' \to B'$ を選べる。
中山の補題（『代数』の補題 \ref{algebra-lemma-NAK}）により、
$B'$ は $P'$ の商である。特に、$A$ 上平坦な $P$ 加群としての $B$ を
持ち上げる、$A'$ 上平坦な $P'$ 加群構造を $B'$ に入れられる。
逆に、$B$ を $A'$ 上平坦な $P'$ 加群 $M'$ に持ち上げられるならば、
$M'$ は巡回加群 $M' \cong P'/J'$ である（ここでも中山の補題を用いる）。
$B' = P'/J'$ とおけば、代数としての $B$ の平坦な持ち上げが得られる。

\medskip\noindent
$C = S^{-1}B$ とおき、$A' \to C'$ を $A \to C$ の平坦な持ち上げとする。
$C$ の可逆元に写る $C'$ の元は可逆である。$A$ 代数準同型
$P \to C$ を持ち上げる $A'$ 代数準同型 $P' \to C'$ を選ぶ。
上の注意により、これらの準同型は局所化を経由し、全射
$S_P^{-1}P \to C$ および $S_{P'}^{-1}P' \to C'$ を与える
（後者には中山の補題を用いる）。特に、$A$ 上平坦な
$S_P^{-1}P$ 加群としての $C$ を持ち上げる、$A'$ 上平坦な
$S_{P'}^{-1}P'$ 加群構造を $C'$ に入れられる。逆に、$C$ を
$A'$ 上平坦な $S_{P'}^{-1}P'$ 加群 $N'$ に持ち上げられるならば、
$N'$ は巡回加群 $N' \cong S_{P'}^{-1}P'/\tilde J$ である
（ここでも中山の補題を用いる）。$C' = S_{P'}^{-1}P'/\tilde J$ とおけば、
代数としての $C$ の平坦な持ち上げが得られる。

\medskip\noindent
スキームの射のシントミック軌跡は定義により開である。
$J_B \subset B$ を、$\Spec(B)$ において
$\Spec(B) \to \Spec(A)$ がシントミックでない点の集合を切り出すイデアルとする。
対応するイデアルを $J_P \subset P$ および $J_{P'} \subset P'$ と書く。
$P' \to S_{P'}^{-1}P'$ は平坦な環準同型であり、補題の $S$ に関する仮定により
同型 $P'/J_{P'} = S_{P'}^{-1}P'/J_{P'}S_{P'}^{-1}P'$ を誘導する。
実際、補題の仮定はちょうど $B/J_B = S^{-1}(B/J_B)$ ということである。
したがって、すべての $g \in J_B$ に対して $B_g$ が $A'$ 上平坦な
$P'_g$ 加群に持ち上がることを示せば、本補題は補題
\ref{examples-defos-lemma-lift-equivalence-module} から従う。
% P11-E0063
しかし環準同型 $A \to B_g$ はシントミックなので、『環準同型の平滑化』の命題
\ref{smoothing-proposition-lift-smooth} により $A'$ 上平坦な代数 $B'$ に持ち上がる。
$A' \to P'_g$ は滑らかなので、前と同様に $P_g \to B_g$ を全射
$P'_g \to B'$ に持ち上げられ、求めるものを得る。
\end{proof}

\begin{lemma}
\label{examples-defos-lemma-first-order-localization}
$k$ を体、$B$ を有限型 $k$ 代数とする。$S \subset B$ を、
$\Spec(B) \to \Spec(k)$ が $\mathfrak q$ において滑らかでなければ
$S \cap \mathfrak q = \emptyset$ となる乗法的部分集合とする。
% P11-E0060
$N$ を有限 $B$ 加群とする。このとき標準的な全単射
$$
\text{Exal}_k(B, N) \to \text{Exal}_k(S^{-1}B, S^{-1}N)
$$
が存在する。
\end{lemma}

\begin{proof}
この証明は補題 \ref{examples-defos-lemma-first-order-completion} の証明と同じだが、より易しい。
読者には証明を飛ばすことを勧める。この写像は局所化によって与えられる。
すなわち、$\text{Exal}_k(B, N)$ の元 $0 \to N \to C \to B \to 0$ を、
$S$ の逆像 $S_C \subset C$ に関する $C$ の局所化 $S_C^{-1}C$ に送る。
補題 \ref{examples-defos-lemma-localization} の証明と比較せよ。

\medskip\noindent
スキームの射の滑らかな軌跡は定義により開である。
$J \subset B$ を、$\Spec(B)$ において
$\Spec(B) \to \Spec(A)$ が滑らかでない点の集合を切り出すイデアルとする。
% P11-E0061
$k \to B$ は有限表示なので、複体 $\NL_{B/k}$ は、$N^i$ が有限 $B$ 加群である
複体 $N^{-1} \to N^0$ によって表せる。『代数』の節
\ref{algebra-section-netherlander}、特に補題
\ref{algebra-lemma-NL-homotopy} を参照せよ。$B$ はネーターなので、
これは $\NL_{B/k}$ が擬連接であることを意味する。$g \in J$ に対して
$k$ 代数 $B_g$ は滑らかであり、したがって
$(\NL_{B/k})_g = \NL_{B_g/k}$ は次数 $0$ に置かれた有限射影 $B$ 加群と
擬同型である。
% P11-E0062
ゆえに任意の $B$ 加群 $N$ と $i \geq 1$ に対して
$\text{Ext}^i_B(\NL_{B/k}, N)_g = 0$ である。最後に次の等式を得る。
$$
\text{Ext}^1_{S^{-1}B}(\NL_{S^{-1}B/k}, S^{-1}N) =
\text{Ext}^1_B(\NL_{B/k}, N) \otimes_B S^{-1}B =
\text{Ext}^1_B(\NL_{B/k}, N)
$$
第1の等式は『代数の詳論』の補題
\ref{more-algebra-lemma-base-change-RHom} と『代数』の補題
\ref{algebra-lemma-localize-NL} による。第2の等式は、
$\text{Ext}^1_B(\NL_{B/k}, N)$ が $J$ 冪捩れであり、$S$ の元が
$J$ 冪捩れ加群上で可逆に作用することによる。補題
\ref{examples-defos-lemma-first-order-completion} の直前で与えた
$\text{Exal}_A(B, N)$ の $\text{Ext}^1_B(\NL_{B/A}, N)$ による記述から、
これで証明が完了する。
\end{proof}

\begin{lemma}
\label{examples-defos-lemma-smooth-localization}
例 \ref{examples-defos-example-rings} において $P$ を $k$ 代数とし、$S \subset P$ を
乗法的部分集合とする。もし
\begin{enumerate}
\item $k \to P$ が有限型であり、
\item すべての $g \in S$ に対して $\Spec(P) \to \Spec(k)$ が
$V(g)$ のすべての点で滑らかである
\end{enumerate}
ならば、補題 \ref{examples-defos-lemma-localization} の変形圏の間の関手
$$
\Deformationcategory_P \longrightarrow \Deformationcategory_{S^{-1}P}
$$
は滑らかであり、接空間上の同型を誘導する。
\end{lemma}

\begin{proof}
補題 \ref{examples-defos-lemma-rings-RS} により、$\Deformationcategory_P$ と
$\Deformationcategory_{S^{-1}P}$ は変形圏である。したがって、
本関手が接空間を同一視し、持ち上げ可能性が対応することを確認すれば十分である。
『形式的変形理論』の補題
\ref{formal-defos-lemma-easy-check-smooth} を参照せよ。持ち上げ可能性に関する性質は
補題 \ref{examples-defos-lemma-lift-equivalence-localization} で示され、接空間上の同型は補題
\ref{examples-defos-lemma-first-order-localization} における $N = B$ の特別な場合である。
\end{proof}



\section{ヘンゼル化の変形}
\label{examples-defos-section-compare-henselization}

\noindent
本節では、代数とその完備化が定める変形問題を比較する。
% P11-E0064
まず「持ち上げ可能性」について論じる。

\begin{lemma}
\label{examples-defos-lemma-lift-equivalence-henselization}
$A' \to A$ を冪零核をもつネーター環の全射とし、$A \to B$ を
有限型の平坦な環準同型とする。$\mathfrak b \subset B$ を、
$\Spec(B) \to \Spec(A)$ が $V(\mathfrak b)$ の補集合上シントミックとなる
イデアルとする。$(B^h, \mathfrak b^h)$ を対 $(B, \mathfrak b)$ のヘンゼル化とする。
このとき $B$ が $A'$ への平坦な持ち上げをもつことと、$B^h$ が
$A'$ への平坦な持ち上げをもつことは同値である。
\end{lemma}

\begin{proof}[第1の証明]
この証明は少しずるい。実際、$B$ が平坦な持ち上げ $B'$ をもつならば、
ヘンゼル化 $(B')^h$ を取ることで $B^h$ の平坦な持ち上げを得る
（補題 \ref{examples-defos-lemma-henselization} の証明と比較せよ）。逆に、$C'$ を
$(B')^h$ の $A'$ 上平坦な持ち上げとする。
% P11-E0065
$\mathfrak c' \subset C'$ をイデアル $\mathfrak b^h$ の逆像とする。
このとき、$\mathfrak c'$ に関する $C'$ の完備化 $(C')^\wedge$ は
$B^\wedge$ の持ち上げである（詳細は省略する）。したがって補題
\ref{examples-defos-lemma-lift-equivalence} により、$B$ は平坦な持ち上げをもつ。
\end{proof}

\begin{proof}[第2の証明]
$A$ 代数の全射 $P = A[x_1, \ldots, x_n] \to B$ を選ぶ。
$\mathfrak p \subset P$ を $\mathfrak b$ の逆像とする。
$P' = A'[x_1, \ldots, x_n]$ とおき、$\mathfrak p' \subset P'$ で
$\mathfrak p$ の逆像を表す（もちろん、ここで $\mathfrak p$ と
$\mathfrak p'$ は素イデアルを表すものではない）。それぞれのヘンゼル化を
$P^h$ および $(P')^h$ と書く。ヘンゼル化を取る操作が関手的であること、
および商のヘンゼル化がヘンゼル化の対応する商であることを用いる。
『代数の詳論』の補題
\ref{more-algebra-lemma-irreducible-henselian-pair-connected} および
\ref{more-algebra-lemma-henselization-integral} を参照せよ。

\medskip\noindent
$A' \to B'$ を $A \to B$ の平坦な持ち上げとする。すなわち、
$A' \to B'$ は平坦で、$A$ 代数の同型
$B = B' \otimes_{A'} A$ が存在するとする。このとき、与えられた全射
$P \to B$ を持ち上げる $A'$ 代数準同型 $P' \to B'$ を選べる。
中山の補題（『代数』の補題 \ref{algebra-lemma-NAK}）により、
$B'$ は $P'$ の商である。特に、$A$ 上平坦な $P$ 加群としての $B$ を
持ち上げる、$A'$ 上平坦な $P'$ 加群構造を $B'$ に入れられる。
逆に、$B$ を $A'$ 上平坦な $P'$ 加群 $M'$ に持ち上げられるならば、
$M'$ は巡回加群 $M' \cong P'/J'$ である（ここでも中山の補題を用いる）。
$B' = P'/J'$ とおけば、代数としての $B$ の平坦な持ち上げが得られる。

\medskip\noindent
$C = B^h$、$\mathfrak c = \mathfrak bC$ とおく。
$A' \to C'$ を $A \to C$ の平坦な持ち上げとする。このとき $C'$ は、
$\mathfrak c$ の逆像 $\mathfrak c'$ に関してヘンゼル的である
（『代数の詳論』の補題
\ref{more-algebra-lemma-henselian-henselian-pair} と、$C' \to C$ の核が
冪零であることによる）。$A$ 代数準同型 $P \to C$ を持ち上げる
$A'$ 代数準同型 $P' \to C'$ を選ぶ。これらの準同型はヘンゼル化を経由し、
全射 $P^h \to C$ および $(P')^h \to C'$ を与える
（後者には再び中山の補題を用いる）。特に、$A$ 上平坦な $P^h$ 加群としての
$C$ を持ち上げる、$A'$ 上平坦な $(P')^h$ 加群構造を $C'$ に入れられる。
逆に、$C$ を $A'$ 上平坦な $(P')^h$ 加群 $N'$ に持ち上げられるならば、
$N'$ は巡回加群 $N' \cong (P')^h/\tilde J$ である
（ここでも中山の補題を用いる）。$C' = (P')^h/\tilde J$ とおけば、
代数としての $C$ の平坦な持ち上げが得られる。

\medskip\noindent
$P' \to (P')^h$ は平坦な環準同型であり、同型
$P'/\mathfrak p' = (P')^h/\mathfrak p'(P')^h$ を誘導する
（『代数の詳論』の補題 \ref{more-algebra-lemma-henselization-flat}）。
したがって、すべての $g \in \mathfrak p'$ に対して $B_g$ が
$A'$ 上平坦な $P'_g$ 加群に持ち上がることを示せば、本補題は補題
\ref{examples-defos-lemma-lift-equivalence-module} から従う。しかし環準同型
$A \to B_g$ はシントミックなので、『環準同型の平滑化』の命題
\ref{smoothing-proposition-lift-smooth} により $A'$ 上平坦な代数 $B'$ に持ち上がる。
$A' \to P'_g$ は滑らかなので、前と同様に $P_g \to B_g$ を全射
$P'_g \to B'$ に持ち上げられ、求めるものを得る。
\end{proof}

\begin{lemma}
\label{examples-defos-lemma-first-order-henselization}
$k$ を体、$B$ を有限型 $k$ 代数とする。$J \subset B$ を、
$\Spec(B) \to \Spec(k)$ が $V(J)$ の補集合上滑らかとなるイデアルとする。
$N$ を有限 $B$ 加群とする。このとき標準的な全単射
$$
\text{Exal}_k(B, N) \to \text{Exal}_k(B^h, N^h)
$$
が存在する。ここで $(B^h, J^h)$ は $(B, J)$ のヘンゼル化であり、
$N^h = N \otimes_B B^h$ である。
\end{lemma}

\begin{proof}
この証明は補題 \ref{examples-defos-lemma-first-order-completion} の証明と同じだが、より易しい。
読者には証明を飛ばすことを勧める。この写像はヘンゼル化によって与えられる。
すなわち、$\text{Exal}_k(B, N)$ の元 $0 \to N \to C \to B \to 0$ を、
$J$ の逆像 $J_C \subset C$ に関する $C$ のヘンゼル化 $C^h$ に送る。
補題 \ref{examples-defos-lemma-henselization} の証明と比較せよ。

\medskip\noindent
$k \to B$ は有限表示なので、複体 $\NL_{B/k}$ は、$N^i$ が有限 $B$ 加群である
複体 $N^{-1} \to N^0$ によって表せる。『代数』の節
\ref{algebra-section-netherlander}、特に補題
\ref{algebra-lemma-NL-homotopy} を参照せよ。$B$ はネーターなので、
これは $\NL_{B/k}$ が擬連接であることを意味する。$g \in J$ に対して
$k$ 代数 $B_g$ は滑らかであり、したがって
$(\NL_{B/k})_g = \NL_{B_g/k}$ は次数 $0$ に置かれた有限射影 $B$ 加群と
擬同型である。
% P11-E0066
ゆえに任意の $B$ 加群 $N$ と $i \geq 1$ に対して
$\text{Ext}^i_B(\NL_{B/k}, N)_g = 0$ である。最後に次の等式を得る。
\begin{align*}
\text{Ext}^1_{B^h}(\NL_{B^h/k}, N^h)
& =
\text{Ext}^1_{B^h}(\NL_{B/k} \otimes_B B^h, N \otimes_B B^h) \\
& =
\text{Ext}^1_B(\NL_{B/k}, N) \otimes_B B^h \\
& =
\text{Ext}^1_B(\NL_{B/k}, N)
\end{align*}
第1の等式は『代数の詳論』の補題
\ref{more-algebra-lemma-henselization-NL}
（より正確には、対のヘンゼル化に対する類似結果）による。第2の等式は
『代数の詳論』の補題 \ref{more-algebra-lemma-base-change-RHom} による。
第3の等式は、$\text{Ext}^1_B(\NL_{B/k}, N)$ が $J$ 冪捩れであること、
写像 $B \to B^h$ が平坦で同型 $B/J \to B^h/JB^h$ を誘導すること
（『代数の詳論』の補題 \ref{more-algebra-lemma-henselization-flat}）、および
『代数の詳論』の補題
\ref{more-algebra-lemma-neighbourhood-equivalence} による。補題
\ref{examples-defos-lemma-first-order-completion} の直前で与えた
$\text{Exal}_A(B, N)$ の $\text{Ext}^1_B(\NL_{B/A}, N)$ による記述から、
これで証明が完了する。
\end{proof}

\begin{lemma}
\label{examples-defos-lemma-smooth-henselization}
例 \ref{examples-defos-example-rings} において $P$ を $k$ 代数とし、$J \subset P$ を
イデアルとする。$(P^h, J^h)$ を対 $(P, J)$ のヘンゼル化とする。もし
\begin{enumerate}
\item $k \to P$ が有限型であり、
\item $\Spec(P) \to \Spec(k)$ が $V(J)$ の補集合上滑らかである
\end{enumerate}
ならば、補題 \ref{examples-defos-lemma-henselization} の変形圏の間の関手
$$
\Deformationcategory_P \longrightarrow \Deformationcategory_{P^h}
$$
は滑らかであり、接空間上の同型を誘導する。
\end{lemma}

\begin{proof}
補題 \ref{examples-defos-lemma-rings-RS} により、$\Deformationcategory_P$ と
$\Deformationcategory_{P^h}$ は変形圏である。したがって、
本関手が接空間を同一視し、持ち上げ可能性が対応することを確認すれば十分である。
『形式的変形理論』の補題
\ref{formal-defos-lemma-easy-check-smooth} を参照せよ。持ち上げ可能性に関する性質は
補題 \ref{examples-defos-lemma-lift-equivalence-henselization} で示され、接空間上の同型は補題
\ref{examples-defos-lemma-first-order-henselization} における $N = B$ の特別な場合である。
\end{proof}



\section{孤立特異点への応用}
\label{examples-defos-section-isolated}

\noindent
上の議論を応用し、有限個の特異点をもつ有限型代数の変形理論を調べる。

\begin{lemma}
\label{examples-defos-lemma-isolated}
例 \ref{examples-defos-example-rings} において $P$ を $k$ 代数とする。
$k \to P$ は有限型であり、$\Spec(P) \to \Spec(k)$ は $P$ の極大イデアル
$\mathfrak m_1, \ldots, \mathfrak m_n$ を除いて滑らかであると仮定する。
$P_{\mathfrak m_i}$、$P_{\mathfrak m_i}^h$、$P_{\mathfrak m_i}^\wedge$ を
それぞれ局所環、ヘンゼル化、完備化とする。このとき変形圏の写像
$$
\Deformationcategory_P \to
\prod \Deformationcategory_{P_{\mathfrak m_i}} \to
\prod \Deformationcategory_{P_{\mathfrak m_i}^h} \to
\prod \Deformationcategory_{P_{\mathfrak m_i}^\wedge}
$$
は滑らかであり、それぞれの有限次元接空間上の同型を誘導する。
\end{lemma}

\begin{proof}
補題 \ref{examples-defos-lemma-finite-type-rings-TI} により接空間は有限次元である。
圏の間の関手は補題 \ref{examples-defos-lemma-localization}、\ref{examples-defos-lemma-henselization}、
\ref{examples-defos-lemma-completion} で構成されている（ヘンゼル化の完備化は完備化である、
という類の確認をいくつか省略する）。

\medskip\noindent
$J = \mathfrak m_1 \cap \ldots \cap \mathfrak m_n$ とおく。補題
\ref{examples-defos-lemma-smooth-completion} を適用すると、
$\Deformationcategory_P \to \Deformationcategory_{P^\wedge}$
は滑らかであり、接空間上の同型を誘導する。ここで $P^\wedge$ は
$P$ の $J$ 進完備化である。しかし $P^\wedge = \prod P_{\mathfrak m_i}^\wedge$
なので、写像 $\Deformationcategory_P \to
\prod \Deformationcategory_{P_{\mathfrak m_i}^\wedge}$
も滑らかであり、接空間上の同型を誘導する。

\medskip\noindent
$(P^h, J^h)$ を対 $(P, J)$ のヘンゼル化とする。このとき
$P^h = \prod P_{\mathfrak m_i}^h$ である（冪等元を調べ、
『代数の詳論』の補題
\ref{more-algebra-lemma-characterize-henselian-pair} を用いよ）。
したがって補題 \ref{examples-defos-lemma-smooth-henselization} を適用し、
完備化の場合と同様に結論できる。

\medskip\noindent
最後の場合を得るには、各 $i$ について個別に
$\Deformationcategory_{P_{\mathfrak m_i}} \to
\Deformationcategory_{P_{\mathfrak m_i}^\wedge}$
が滑らかであり、接空間上の同型を誘導することを示せば十分である。
% P11-E0067
そのため、$P$ を、唯一の特異点が極大イデアル $\mathfrak m$ であるような
単項局所化に置き換えてよい（これは元の $P$ における $\mathfrak m_i$ に対応する）。
そこで、乗法的部分集合 $S = P \setminus \mathfrak m$ に補題
\ref{examples-defos-lemma-smooth-localization} を適用すれば結論を得る。細部は省略する。
\end{proof}






\section{障害のない変形問題}
\label{examples-defos-section-unobstructed}

\noindent
$p : \mathcal{F} \to \mathcal{C}_\Lambda$ を、グルーポイドをファイバーとする
余ファイバー圏とする。$p$ が滑らかなとき、$\mathcal{F}$ は
{\it 滑らか}または{\it 障害がない}ということを思い出そう。これは、
$\mathcal{C}_\Lambda$ における全射 $\varphi : A' \to A$ と
$x \in \Ob(\mathcal{F}(A))$ が与えられたとき、$p(f) = \varphi$ を満たす
$\mathcal{F}$ の射 $f : x' \to x$ が存在することを意味する。
『形式的変形理論』の節 \ref{formal-defos-section-smooth} を参照せよ。
本節では、幾何学的に意味のある例をいくつか与える。

\begin{lemma}
\label{examples-defos-lemma-lci-unobstructed}
例 \ref{examples-defos-example-rings} において、$P$ を $k$ 上の局所完全交叉とする
（『代数』の定義 \ref{algebra-definition-lci-field}）。このとき
$\Deformationcategory_P$ には障害がない。
\end{lemma}

\begin{proof}
$(A, Q) \to (k, P)$ を $\Deformationcategory_P$ の対象とする。
『代数』の定義 \ref{algebra-definition-lci} により、$A \to Q$ は
シントミックな環準同型である。したがって、$\mathcal{C}_\Lambda$ の任意の全射
$A' \to A$ に対し、『環準同型の平滑化』の命題
\ref{smoothing-proposition-lift-smooth} により、$A' \to A$ を持ち上げる射
$(A', Q') \to (A, Q)$ が存在する。これで補題が証明された。
\end{proof}

\begin{lemma}
\label{examples-defos-lemma-glueing-smooth}
状況 \ref{examples-defos-situation-glueing} において $U_{12} \to \Spec(k)$ が滑らかならば、射
$$
\Deformationcategory_X
\longrightarrow
\Deformationcategory_{U_1} \times \Deformationcategory_{U_2} =
\Deformationcategory_{P_1} \times \Deformationcategory_{P_2}
$$
は滑らかである。さらに $U_1$ が $k$ 上の局所完全交叉ならば、
$$
\Deformationcategory_X
\longrightarrow
\Deformationcategory_{U_2} = \Deformationcategory_{P_2}
$$
は滑らかである。
\end{lemma}

\begin{proof}
等号は補題 \ref{examples-defos-lemma-affine} により成り立つ。『形式的変形理論』の節
\ref{formal-defos-section-smooth} のように、$\mathcal{C}_\Lambda$ を
$\mathcal{C}_\Lambda$ 上の変形圏とみなす。このとき
$$
\Deformationcategory_{P_1} \times \Deformationcategory_{P_2} =
\Deformationcategory_{P_1}
\times_{\mathcal{C}_\Lambda}
\Deformationcategory_{P_2},
$$
である。『形式的変形理論』の注意
\ref{formal-defos-remarks-cofibered-groupoids}
（\ref{formal-defos-item-product}）を参照せよ。補題
\ref{examples-defos-lemma-glueing} を用いると、第1の主張は、関手
$$
\Deformationcategory_{P_1}
\times_{\Deformationcategory_{P_{12}}}
\Deformationcategory_{P_2}
\longrightarrow
\Deformationcategory_{P_1}
\times_{\mathcal{C}_\Lambda}
\Deformationcategory_{P_2}
$$
が滑らかであるということである。これは
$T\Deformationcategory_{P_{12}} = (0)$ を示せば、『形式的変形理論』の補題
\ref{formal-defos-lemma-map-fibre-products-smooth} から従う。この消滅は、
$P_{12}$ が $k$ 上滑らかなので補題 \ref{examples-defos-lemma-smooth} から従う。
第2の主張については、
$\Deformationcategory_{P_1} \to \mathcal{C}_\Lambda$
が滑らかであることを示せば十分である。『形式的変形理論』の補題
\ref{formal-defos-lemma-smooth-properties} を参照せよ。言い換えれば、
$\Deformationcategory_{P_1}$ に障害がないことを示す必要があるが、
これは補題 \ref{examples-defos-lemma-lci-unobstructed} である。
\end{proof}

\begin{lemma}
\label{examples-defos-lemma-curve-isolated}
例 \ref{examples-defos-example-schemes} において $X$ を $k$ 上のスキームとする。次を仮定する。
\begin{enumerate}
\item $X$ は分離的かつ $k$ 上有限型で、$\dim(X) \leq 1$ である。
\item $X \to \Spec(k)$ は閉点 $p_1, \ldots, p_n \in X$ を除いて滑らかである。
\end{enumerate}
$\mathcal{O}_{X, p_1}$、$\mathcal{O}_{X, p_1}^h$、
$\mathcal{O}_{X, p_1}^\wedge$ を、それぞれ局所環、ヘンゼル化、完備化とする。
% P11-E0068
変形圏の写像
$$
\Deformationcategory_X
\longrightarrow
\prod \Deformationcategory_{\mathcal{O}_{X, p_i}}
\longrightarrow
\prod \Deformationcategory_{\mathcal{O}_{X, p_i}^h}
\longrightarrow
\prod \Deformationcategory_{\mathcal{O}_{X, p_i}^\wedge}
$$
を考える。第1の射は滑らかであり、第2と第3の射は滑らかで、
接空間上の同型を誘導する。
\end{lemma}

\begin{proof}
$p_1, \ldots, p_n$ と $X$ の各既約成分の生成点を含むアフィン開部分
$U_2 \subset X$ を選ぶ。これは『多様体』の補題
\ref{varieties-lemma-dim-1-quasi-projective} と『性質』の補題
\ref{properties-lemma-ample-finite-set-in-affine} により可能である。
このとき $X \setminus U_2$ は有限であり、$X = U_1 \cup U_2$ となる
アフィン開部分 $U_1 \subset X \setminus \{p_1, \ldots, p_n\}$ を選べる。
$U_{12} = U_1 \cap U_2$ とおく。このとき $U_1$ と $U_{12}$ は
$k$ 上の滑らかなアフィンスキームである。補題 \ref{examples-defos-lemma-glueing-smooth} により
$$
\Deformationcategory_X \longrightarrow \Deformationcategory_{U_2}
$$
は滑らかである。補題 \ref{examples-defos-lemma-affine} と \ref{examples-defos-lemma-isolated} を
適用すれば結論を得る。
\end{proof}

\begin{lemma}
\label{examples-defos-lemma-curve-isolated-lci}
例 \ref{examples-defos-example-schemes} において $X$ を $k$ 上のスキームとする。次を仮定する。
\begin{enumerate}
\item $X$ は分離的かつ $k$ 上有限型で、$\dim(X) \leq 1$ である。
\item $X$ は $k$ 上の局所完全交叉である。
\item $X \to \Spec(k)$ は有限個の点を除いて滑らかである。
\end{enumerate}
このとき $\Deformationcategory_X$ には障害がない。
\end{lemma}

\begin{proof}
$p_1, \ldots, p_n \in X$ を $X \to \Spec(k)$ が滑らかでない点とする。
$p_1, \ldots, p_n$ と $X$ の各既約成分の生成点を含むアフィン開部分
$U_2 \subset X$ を選ぶ。これは『多様体』の補題
\ref{varieties-lemma-dim-1-quasi-projective} と『性質』の補題
\ref{properties-lemma-ample-finite-set-in-affine} により可能である。
このとき $X \setminus U_2$ は有限であり、$X = U_1 \cup U_2$ となる
アフィン開部分 $U_1 \subset X \setminus \{p_1, \ldots, p_n\}$ を選べる。
$U_{12} = U_1 \cap U_2$ とおく。このとき $U_1$ と $U_{12}$ は
$k$ 上の滑らかなアフィンスキームである。補題 \ref{examples-defos-lemma-glueing-smooth} により
$$
\Deformationcategory_X \longrightarrow \Deformationcategory_{U_2}
$$
は滑らかである。補題 \ref{examples-defos-lemma-affine} と \ref{examples-defos-lemma-lci-unobstructed} を
適用すれば結論を得る。
\end{proof}




\section{平滑化}
\label{examples-defos-section-smoothing}

\noindent
体 $k$ 上の有限型スキームまたは代数空間 $X$ が与えられたとする。
一般ファイバーが滑らかで、特殊ファイバーが $X$ と同型である有限型平坦射
$Y \to \Spec(k[[t]])$ を見つけることがしばしば有用である。このようなものを
$X$ の平滑化という。本節では、孤立した局所完全交叉特異点をもつ
$1$ 次元分離的な $X$ の平滑化を見つける。

\begin{lemma}
\label{examples-defos-lemma-criterion-smoothing}
$k$ を体とし、$S = \Spec(k[[t]])$、$S_n = \Spec(k[t]/(t^n))$ とおく。
$Y \to S$ をスキームの固有平坦射とし、その特殊ファイバー $X$ は
Cohen--Macaulay かつ次元 $d$ の等次元であるとする。
$X_n = Y \times_S S_n$ と書く。ある $n \geq 1$ に対して
$\Omega_{X_n/S_n}$ の第 $d$ Fitting イデアルが $t^{n - 1}$ を含むならば、
$Y \to S$ の一般ファイバーは滑らかである。
\end{lemma}

\begin{proof}
『射の詳論』の補題
\ref{more-morphisms-lemma-flat-finite-presentation-CM-open}
により、$Y \to S$ は Cohen--Macaulay 射である。『射』の補題
\ref{morphisms-lemma-flat-finite-presentation-CM-fibres-relative-dimension}
により、$Y \to S$ は相対次元 $d$ である。『因子』の補題
\ref{divisors-lemma-d-fitting-ideal-omega-smooth} により、
$\Omega_{Y/S}$ の第 $d$ Fitting イデアル
$\mathcal{I} \subset \mathcal{O}_Y$ は、射 $Y \to S$ の特異軌跡を切り出す。
言い換えれば、$V(\mathcal{I}) \subset Y$ は $Y \to S$ が滑らかでない点からなる
閉部分集合である。『因子』の補題
\ref{divisors-lemma-base-change-and-fitting-ideal-omega} により、この Fitting イデアルの
形成は基底変換と可換である。仮定から、$t^{n - 1}$ は
$\mathcal{I} + t^n\mathcal{O}_Y$ の切断である。したがって、すべての
$x \in X = V(t) \subset Y$ に対し、$\mathcal{I}_x$ を $x$ における茎とすると
$t^{n - 1} \in \mathcal{I}_x$ である。これは $Y$ における $X$ のある開近傍で
$V(\mathcal{I}) \subset V(t)$ となることを意味する。$Y \to S$ は固有なので、
求める $V(\mathcal{I}) \subset V(t)$ が従う。
\end{proof}

\begin{lemma}
\label{examples-defos-lemma-jouanolou-type-thing}
$k$ を体とし、$1 \leq c \leq n$ を整数とする。
$f_1, \ldots, f_c \in k[x_1, \ldots x_n]$ を元とする。
% P11-E0069
$a_{ij}$（$0 \leq i \leq n$、$1 \leq j \leq c$）を変数とする。次を考える。
$$
g_j = f_j + a_{0j} + a_{1j}x_1 + \ldots + a_{nj}x_n \in
k[a_{ij}][x_1, \ldots, x_n]
$$
$Y \subset \mathbf{A}^{n + c(n + 1)}_k$ を $g_1, \ldots, g_c$ が切り出す
閉部分スキームとする。$\pi : Y \to \mathbf{A}^{c(n + 1)}_k$ を、
変数 $a_{ij}$ をもつアフィン空間への射影とする。このとき
$\mathbf{A}^{c(n + 1)}_k$ の空でない Zariski 開部分集合で、
$\pi$ が滑らかとなるものが存在する。
% P11-E0070
\end{lemma}

\begin{proof}
$\pi$ が滑らかである点の集合は開であることを思い出そう。したがってその補集合、
すなわち特異軌跡は閉である。Chevalley の定理
（『射』の補題 \ref{morphisms-lemma-chevalley} の形）により、特異軌跡の像は
構成可能である。したがって、$\mathbf{A}^{c(n + 1)}_k$ の生成点が特異軌跡の像に
属さなければ、補題が従う（例えば『位相』の補題
\ref{topology-lemma-generic-point-in-constructible} による）。ゆえに、
$\pi$ が滑らかでなく、$\mathbf{A}^{c(n + 1)}_k$ の生成点に写る点
$y \in Y$ が存在しないことを示す必要がある。偏微分の行列
$$
(\frac{\partial g_j}{\partial x_i}) =
(\frac{\partial f_j}{\partial x_i} + a_{ij})
$$
$\kappa(y)$ におけるこの行列の像の階数は $< c$ でなければならない。
そうでなければ $\pi$ は $y$ において滑らかだからである。
『環準同型の平滑化』の節 \ref{smoothing-section-singular-ideal} の議論を参照せよ。
したがって、すべてが零ではない $\lambda_1, \ldots, \lambda_c \in \kappa(y)$ で、
ベクトル $(\lambda_1, \ldots, \lambda_c)$ がこの行列の核に属するものを選べる。
番号を付け替えて $\lambda_1 \not = 0$ と仮定してよい。
$\lambda_1$ で割ることにより、ベクトルが
$(1, \lambda_2, \ldots, \lambda_c)$ の形であると仮定してよい。このとき
$$
a_{i1} = -
\frac{\partial f_j}{\partial x_1} -
\sum\nolimits_{j = 2, \ldots, c} \lambda_j(\frac{\partial f_j}{\partial x_i} + a_{ij})
$$
を $\kappa(y)$ において $i = 1, \ldots, n$ に対して得る。
% P11-E0071
さらに $y \in Y$ なので、
$$
a_{0j} = -f_j - a_{1j}x_1 - \ldots - a_{nj}x_n
$$
を $\kappa(y)$ において得る。これは次を意味する。

$a_{ij}$ が生成する $\kappa(y)$ の部分体は、
$x_1, \ldots, x_n, \lambda_2, \ldots, \lambda_c$ の像と、$a_{i1}$ および
$a_{0j}$ を除く $a_{ij}$ が生成する $\kappa(y)$ の部分体に含まれることを意味する。
個数を数えると、この超越次数は高々 $c(n + 1) - 1$ である。したがって
$y$ は生成点に写ることができず、求める結論を得る。
\end{proof}

\begin{lemma}
\label{examples-defos-lemma-smoothing-affine-lci}
$k$ を体とし、$A$ を $k$ 上の大域的完全交叉とする。
$B/tB \cong A$ を満たす有限型平坦な環準同型 $k[[t]] \to B$ で、
$B[1/t]$ が $k((t))$ 上滑らかとなるものが存在する。
\end{lemma}

\begin{proof}
『代数』の定義 \ref{algebra-definition-lci-field} のように
$A = k[x_1, \ldots, x_n]/(f_1, \ldots, f_c)$ と書く。
$a_{ij} \in (t) \subset k[[t]]$ を選び、
$$
g_j = f_j + a_{0j} + a_{1j}x_1 + \ldots + a_{nj}x_n \in
k[[t]][x_1, \ldots, x_n]
$$
とおく。その後、$B = k[[t]][x_1, \ldots, x_n]/(g_1, \ldots, g_c)$ とする。
$k[[t]] \to B$ は $(t)$ の上にあるすべての素イデアルにおいて平坦であると主張する。
実際、$f_1, \ldots, f_c$ を含む $k[x_1, \ldots, x_n]$ の任意の素イデアル
$\mathfrak p$ における局所環で、$f_1, \ldots, f_c$ は正則列をなす
（『代数』の補題 \ref{algebra-lemma-lci}）。したがって
$g_1, \ldots, g_c$ は局所的に正則列の持ち上げであり、『代数』の補題
\ref{algebra-lemma-grothendieck-regular-sequence} を適用できる。
$(0) \subset k[[t]]$ の上にある素イデアルでの平坦性は、
$k((t)) = k[[t]]_{(0)}$ が体なので自動的である。したがって $B$ は
$k[[t]]$ 上平坦である。

\medskip\noindent
残るのは、$a_{ij}$ を適切に選べば一般ファイバー $B_{(0)}$ が
$k((t))$ 上滑らかになることを示すことだけである。そのためには、誘導される射
$$
(a_{ij}) : \Spec(k[[t]]) \longrightarrow \mathbf{A}^{c(n + 1)}_k
$$
が補題 \ref{examples-defos-lemma-jouanolou-type-thing} の空でない Zariski 開部分集合に
入るよう $a_{ij}$ を選べることを示す必要がある。これは、$a_{ij}$ の
非零多項式で $(t)^{\oplus c(n + 1)}$ 上消えるものは存在しないので明らかである
（読者への演習とする）。
\end{proof}

\begin{lemma}
\label{examples-defos-lemma-smoothing-artinian-lci}
$k$ を体とし、$A$ を $k$ 上の局所完全交叉である有限次元 $k$ 代数とする。
このとき $B/tB \cong A$ を満たす有限平坦 $k[[t]]$ 代数 $B$ で、
$B[1/t]$ が $k((t))$ 上エタールとなるものが存在する。
\end{lemma}

\begin{proof}
$A$ はアルティンである（『代数』の補題
\ref{algebra-lemma-finite-dimensional-algebra}）から、$A$ を局所アルティン環の積として
書ける（『代数』の補題 \ref{algebra-lemma-artinian-finite-length}）。
したがって $A$ が局所的な場合に補題を証明すれば十分である
（ここでは、単項局所化を取っても局所完全交叉であることが保たれることを用いる。
『代数』の補題 \ref{algebra-lemma-localize-lci} を参照せよ）。この場合、$A$ は
大域的完全交叉である。補題 \ref{examples-defos-lemma-smoothing-affine-lci} で構成した代数 $B$ を
考える。このとき $k[[t]] \to B$ は、$(t)$ の上にある $B$ の唯一の素イデアルで
準有限である（『代数』の定義 \ref{algebra-definition-quasi-finite}）。
$k[[t]]$ はヘンゼル局所環であることに注意する
（『代数』の補題 \ref{algebra-lemma-complete-henselian}）。したがって
$B = B' \times C$ と書け、$B'$ は $k[[t]]$ 上有限で、$C$ は $(t)$ の上に
素イデアルをもたない。『代数』の補題
\ref{algebra-lemma-characterize-henselian} を参照せよ。この $B'$ が求める環である
（エタールであることは相対次元 $0$ で滑らかであることと同じだと思い出そう）。
\end{proof}

\begin{lemma}
\label{examples-defos-lemma-smoothing-at-lci-point}
$k$ を体とし、$A$ を $k$ 代数とする。次を仮定する。
\begin{enumerate}
\item $A$ は本質的有限型な $k$ 上の局所環である。
\item $A$ は $k$ 上の完全交叉である
（『代数』の定義 \ref{algebra-definition-lci-local-ring}）。
\end{enumerate}
$\kappa$ を $A$ の剰余体とし、$d = \dim(A) + \text{trdeg}_k(\kappa)$ とおく。
このとき整数 $n$ と、本質的有限型かつ平坦な環準同型
$k[[t]] \to B$ で $B/tB \cong A$ を満たし、$t^n$ が
$\Omega_{B/k[[t]]}$ の第 $d$ Fitting イデアルに属するものが存在する。
% P11-E0072
\end{lemma}

\begin{proof}
『代数』の補題 \ref{algebra-lemma-lci-local} により、$A$ を $k$ 上の
大域的完全交叉 $P$ の素イデアル $\mathfrak p$ における局所化として書ける。
『代数』の補題 \ref{algebra-lemma-dimension-at-a-point-finite-type-field} により
$\dim(P) = d$ であることに注意する。補題
\ref{examples-defos-lemma-smoothing-affine-lci} により、$P \cong Q/tQ$ を満たし
$k((t)) \to Q[1/t]$ が滑らかとなる有限型平坦な環準同型
$k[[t]] \to Q$ を見つけられる。同補題における $Q$ の構成から、
$k[[t]] \to Q$ は相対次元 $d$ の相対大域的完全交叉である。別の議論として、
『代数』の補題 \ref{algebra-lemma-syntomic} によれば、$Q$ または $Q$ の
適切な単項局所化がそのような大域的完全交叉である。したがって『因子』の補題
\ref{divisors-lemma-d-fitting-ideal-omega-smooth} により、
$\Omega_{Q/k[[t]]}$ の第 $d$ Fitting イデアル $I \subset Q$ は、
$\Spec(Q) \to \Spec(k[[t]])$ の特異軌跡を切り出す。ゆえに、ある $n$ に対して
$t^n \in I$ である。$\mathfrak q \subset Q$ を $\mathfrak p$ の逆像とし、
$B = Q_\mathfrak q$ とおく。これで補題が証明された。
\end{proof}

\begin{lemma}
\label{examples-defos-lemma-smoothing-proper-curve-isolated-lci}
体 $k$ 上のスキーム $X$ について、次を仮定する。
\begin{enumerate}
\item $X$ は $k$ 上固有である。
\item $X$ は $k$ 上の局所完全交叉である。
\item $X$ の次元は $\leq 1$ である。
\item $X \to \Spec(k)$ は有限個の点を除いて滑らかである。
\end{enumerate}
このとき、一般ファイバーが滑らかで特殊ファイバーが $X$ と同型である
射影的平坦射 $Y \to \Spec(k[[t]])$ が存在する。
\end{lemma}

\begin{proof}
$X$ は Cohen--Macaulay であることに注意する。『代数』の補題
\ref{algebra-lemma-lci-CM} を参照せよ。したがって
$X = X' \amalg X''$ と書け、$\dim(X') = 0$ かつ $X''$ は次元 $1$ の
等次元である。『射』の補題
\ref{morphisms-lemma-flat-finite-presentation-CM-fibres-relative-dimension} を参照せよ。
$X'$ は $k$ 上有限なので（『多様体』の補題
\ref{varieties-lemma-algebraic-scheme-dim-0}）、補題
\ref{examples-defos-lemma-smoothing-artinian-lci} により、特殊ファイバーが $X'$ で一般ファイバーが
滑らかな $Y' \to \Spec(k[[t]])$ を見つけられる。したがって $X''$ について
補題を証明すれば十分である。$X$ を $X''$ に置き換えると、$X$ は
Cohen--Macaulay かつ次元 $1$ の等次元となる。

\medskip\noindent
$\Lambda = k \to k$ という状況に対する変形理論を用いる。
$p_1, \ldots, p_r \in X$ を $X$ の閉特異点、すなわち
$X \to \Spec(k)$ が滑らかでない点とする。各 $i$ に対し、整数 $n_i$ と
本質的有限型な平坦環準同型
$$
k[[t]] \longrightarrow B_i
$$
で、$B_i/tB_i \cong \mathcal{O}_{X, p_i}$ を満たし、$t^{n_i}$ が
$\Omega_{B_i/k[[t]]}$ の第 $1$ Fitting イデアルに属するものを選ぶ。
これは補題 \ref{examples-defos-lemma-smoothing-at-lci-point} により可能である。
系 $(B_i/t^nB_i)$ は $k[[t]]$ 上の
$\Deformationcategory_{\mathcal{O}_{X, p_i}}$ の形式対象を定めることに注意する。
補題 \ref{examples-defos-lemma-curve-isolated} により写像
$$
\Deformationcategory_X
\longrightarrow
\prod\nolimits_{i = 1, \ldots, r} \Deformationcategory_{\mathcal{O}_{X, p_i}}
$$
は変形圏の間の滑らかな写像である。したがって『形式的変形理論』の補題
\ref{formal-defos-lemma-smooth-morphism-essentially-surjective}
により、$\Deformationcategory_X$ の形式対象 $(X_n)$ で、上の射によって形式対象
$\prod_i (B_i/t^n)$ に写るものが存在する。『空間の射の詳論』の補題
\ref{spaces-more-morphisms-lemma-formal-algebraic-space-proper-reldim-1}
により、$k[[t]]$ 上の射影スキーム $Y$ と、両立する同型
$Y \times_{\Spec(k[[t]])} \Spec(k[t]/(t^n)) \cong X_n$ が存在する。
『射の詳論』の補題
\ref{more-morphisms-lemma-check-flatness-on-infinitesimal-nbhds}
により、$Y \to \Spec(k[[t]])$ は平坦である。
$X$ は Cohen--Macaulay かつ次元 $1$ の等次元なので、補題
\ref{examples-defos-lemma-criterion-smoothing} を適用し、$Y$ の一般ファイバーが滑らかであることを
確認できる\footnote{注意：一般には、点 $p_i$ における $Y$ の局所環が
$B_i$ と同型であるとは{\bf 限らない}。分かっているのは、両辺を $t^n$ で
割った後に同型となることだけである。}。上で得た整数 $n_i$ の最大値より
真に大きい $n$ を選ぶ。$S_n = \Spec(k[t]/(t^n))$ としたとき、
$t^{n - 1}$ が $\Omega_{X_n/S_n}$ の第1 Fitting イデアルに属することを
示せれば、証明は完了する。そのためには、$X_n$ の各閉点 $p$ における局所環で
この主張を示せば十分である。一方、$p$ が $X \to \Spec(k)$ の滑らかな点に
対応するならば、$\Omega_{X_n/S_n, p}$ は階数 $1$ の自由加群で、第1 Fitting
イデアルは局所環全体である。ある $i$ に対して $p = p_i$ ならば、
$$
\Omega_{X_n/S_n, p_i} =
\Omega_{(B_i/t^nB_i)/(k[t]/(t^n))} =
\Omega_{B_i/k[[t]]}/t^n\Omega_{B_i/k[[t]]}
$$
Fitting イデアルを取る操作は基底変換と可換であり
（これはすでに用いたが、この代数的な設定では『代数の詳論』の補題
\ref{more-algebra-lemma-fitting-ideal-basics} から従う）、
% P11-E0073
$n - 1 \geq n_i$ なので、$t^{n - 1}$ は求めるとおり
$B_i/t^nB_i$ 上のこの加群の Fitting イデアルに属する。
\end{proof}

\begin{lemma}
\label{examples-defos-lemma-smoothing-curve-isolated-lci}
$k$ を体とし、$X$ を $k$ 上のスキームとする。次を仮定する。
\begin{enumerate}
\item $X$ は分離的かつ $k$ 上有限型で、$\dim(X) \leq 1$ である。
\item $X$ は $k$ 上の局所完全交叉である。
\item $X \to \Spec(k)$ は有限個の点を除いて滑らかである。
\end{enumerate}
このとき、一般ファイバーが滑らかで特殊ファイバーが $X$ と同型である、
分離的な有限型平坦射 $Y \to \Spec(k[[t]])$ が存在する。
\end{lemma}

\begin{proof}
$X$ が被約ならば、『多様体』の補題
\ref{varieties-lemma-reduced-dim-1-projective-completion} のような埋め込み
$X \subset \overline{X}$ を選べる。$X = \overline{X} \setminus \{x_1, \ldots, x_n\}$
と書くと、$\mathcal{O}_{\overline{X}, x_i}$ は離散付値環であり、したがって特に
局所完全交叉である（『代数』の定義
\ref{algebra-definition-lci-local-ring}）。ゆえに $\overline{X}$ は $k$ 上の
局所完全交叉である。実際、開部分 $X$ 上では仮定によりそうであり、点 $x_i$ では
『代数』の補題 \ref{algebra-lemma-lci-local} による。したがって補題
\ref{examples-defos-lemma-smoothing-proper-curve-isolated-lci} を適用し、一般ファイバーが滑らかで
特殊ファイバーが $\overline{X}$ である射影的平坦射
$\overline{Y} \to \Spec(k[[t]])$ を得る。次に $\overline{Y}$ から
$x_1, \ldots, x_n$ を除いて $Y$ を得る。

\medskip\noindent
一般の場合、$X = X' \amalg X''$ と書き、$\dim(X') = 0$ かつ $X''$ は
次元 $1$ の等次元とする。
% P11-E0074
このとき $X''$ は被約であり、第1段落を適用できる。一方、$X'$ は補題
\ref{examples-defos-lemma-smoothing-proper-curve-isolated-lci} の証明と同様に扱える。
細部は省略する。
\end{proof}
